\documentclass[11pt]{article}

\usepackage[a4paper,margin=1in]{geometry}
\usepackage{amsmath,amssymb,amsthm,mathtools}
\usepackage[hidelinks]{hyperref}
\hypersetup{
  pdfauthor={Bangyen Pham},
  pdftitle={Sharpness of the Coefficient Order-Statistic Bound},
  pdfsubject={Sharpness and necessity results for the coefficient order-statistic bound on polynomial multiples},
  pdfkeywords={polynomial multiples, coefficient mass, reduced height, Chebyshev systems, exponential sums},
  pdflang={en-US}
}

\newtheorem{theorem}{Theorem}[section]
\newtheorem{lemma}[theorem]{Lemma}
\newtheorem{corollary}[theorem]{Corollary}
\newtheorem{proposition}[theorem]{Proposition}
\theoremstyle{definition}
\newtheorem{conjecture}[theorem]{Conjecture}
\theoremstyle{plain}
\newcommand{\R}{\mathbb R}
\newcommand{\N}{\mathbb N}
\newcommand{\E}{\mathbb E}
\DeclareMathOperator{\Tail}{Tail}
\DeclareMathOperator{\sgn}{sgn}

\title{Sharpness of the Coefficient Order-Statistic Bound}
\author{\shortstack{Bangyen Pham\\
\small Independent Researcher\\
\small\href{mailto:bangyenp@gmail.com}{\texttt{bangyenp@gmail.com}}\\
\small\href{https://orcid.org/0009-0006-9928-2088}{ORCID: 0009-0006-9928-2088}}}
\date{September 2026}

\begin{document}
\maketitle

\begin{abstract}
For real roots $2\le r_1\le\cdots\le r_L$ and a nonzero multiple $F$ of
$\prod_{i=1}^L(x-r_i)$ with leading coefficient $f_D$, the companion
paper~\cite{coefficient-mass} proves that the $(u+1)$-st largest nonleading
coefficient magnitude satisfies $b_{u+1}\ge|f_D|\prod_{i>u}(r_i-1)$.  We show
that this bound is the infimum over monic multiples exactly when no partial
sum of the coefficients of the surviving factor $P_u=\prod_{i>u}(x-r_i)$
exceeds their total, repeated roots included.  When this criterion fails, the placement in which the exempt coefficients escape
to the bottom is always extremal, so the infimum is the reciprocal of the
least $\ell^1$ tail of a certificate on the surviving roots alone,
equivalently the least, over monic multipliers $M$, of the largest partial
sum of the coefficients of $P_uM$; at $(2,3,5,7)$ it is $1704/31$.  Below $2$
the bound can fail once $u\ge1$, and the node cutoff $1/2$ of the underlying tail comparison cannot
be raised.  For every $r_1>1$ we prove a repaired bound and a sharper
two-sided bound $b_{u+1}\ge|f_D|\min(|P(1)|,|P_u(1)|)$, and determine
exactly when each is the infimum, repeated roots included.  We also extend the tail
comparison to confluent nodes, with strictness.
\end{abstract}

\noindent\textbf{Keywords.} Polynomial multiples; coefficient mass; reduced
height; Chebyshev systems; exponential sums.

\noindent\textbf{2020 Mathematics Subject Classification.} Primary 11C08; Secondary 26C05, 41A50.

\section{Introduction}

The certificate paper~\cite{coefficient-mass} proves a degree-independent
lower bound on the coefficient order statistics of a polynomial multiple.
For real roots $2\le r_1\le\cdots\le r_L$, repetitions allowed, and $F$ a
nonzero multiple of $\prod_{i=1}^L(x-r_i)$ of degree $D\ge L$ with leading
coefficient $f_D$, its Theorem~3.2 gives, for every $0\le u\le L-1$,
\begin{equation}\label{eq:order}
  b_{u+1}\ge|f_D|\prod_{i=u+1}^L(r_i-1),
\end{equation}
where $b_k$ is the $k$th largest of the nonleading coefficient magnitudes
$|f_j|$, $j<D$, so that \eqref{eq:order} bounds $b_k$ with $k=u+1$; the case
$u=0$ needs only $r_1>1$.  The engine behind
\eqref{eq:order} is its Theorem~2.2, a comparison between the infinite
$\ell^1$ tail of a normalized exponential sum with prescribed zeros and the
tail of the sum with the same zero count but a consecutive zero set: for
nodes $0<y_1<\cdots<y_c\le1/2$, $y_i=\rho_i^{-1}$, and a zero set $Z$ of size
$c-1$ with leading run $f=\max\{0\le k\le c-1:\{1,\ldots,k\}\subseteq Z\}$,
the unique exponential sum $u$ on those nodes with $u_0=1$ and $u_z=0$
($z\in Z$) satisfies
\[
  \Tail(u)=\sum_{d\ge1}|u_d|\le\prod_{i=1}^{f+1}\frac{y_i}{1-y_i},
\]
with equality exactly at $Z=\{1,\ldots,c-1\}$; \eqref{eq:order} follows by
applying this with the excluded coefficient positions supplying $Z$.
Lemma~2.1 evaluates the consecutive case; Lemma~2.3 supplies the interlacing
step inside the comparison; Lemma~3.1 converts the tail bound into a
coefficient bound; Corollary~3.6 shows the resulting lower bound gives the
infimum of the logarithmic coefficient mass over monic multiples of the first
$L$ primes up to $O(L^2)$.  Throughout, an unqualified Lemma~2.1,
Theorem~2.2, Lemma~2.3, Lemma~3.1, Theorem~3.2 or Corollary~3.6 is the
result of that number in~\cite{coefficient-mass}; every other numbered result
is this paper's own.

\paragraph{Main results.}  Write $P=\prod_{i=1}^L(x-r_i)$, let
$P_u=\prod_{i>u}(x-r_i)$ be the factor of the $m=L-u$ surviving roots, and let
$S_j(P_u)$ be the partial sums of its coefficients read down from the leading
one, so that $S_j(P_u)=P_u(1)=(-1)^m\prod_{i>u}(r_i-1)$ once $j\ge m$.  All
infima below are of $b_{u+1}(QP)$ over monic $Q\in\R[x]$.

\emph{The criterion.}  For $2\le r_1\le\cdots\le r_L$, repetitions allowed,
\eqref{eq:order} is the infimum if and only if
\[
  |S_j(P_u)|\le|P_u(1)|\qquad(1\le j<m),
\]
which is condition \eqref{eq:partial} below (Proposition~\ref{prop:partial},
Theorem~\ref{thm:converse} and Corollary~\ref{cor:charrep}).  The multiples
approaching the bound carry a short prefix $S_1(P_u),\ldots,S_{m-1}(P_u)$, a
long run of the value $P_u(1)$, and the $u$ exempt coefficients at the
bottom; at $(2,3,5)$ they are correctly ordered at every degree
(Proposition~\ref{prop:sharp235}).  The criterion holds at $(3,4,5,6)$ and
fails at $(2,3,5,7)$ with $u=1$.

\emph{The value.}  For $2\le r_1\le\cdots\le r_L$ let $\tau^*$ be the
least tail $\sum_{d\ge1}|v_d|$ of an exponential sum
$v_d=\sum_{i>u}\lambda_ir_i^{-d}$ with $v_0=1$ (for repeated roots, of an
element of their Hermite space).  Then
\[
  \inf_Qb_{u+1}(QP)=\frac1{\tau^*}=\min_M\max_{j\ge1}\bigl|S_j(P_uM)\bigr|
\]
over monic $M\in\R[x]$ (Theorems~\ref{thm:value} and~\ref{thm:extremal},
Proposition~\ref{prop:attain}, and Theorem~\ref{thm:extremalrep} at repeated
roots): the worst placement of the $u$ largest
coefficients is the one in which they escape to the bottom, and the infimum
depends on the surviving roots alone.  At $(2,3,5,7)$ with $u=1$ it is
$1704/31$ against the bound $48$, and at $(2,\ldots,7)$ with $u=2$ it is
$114840/259$ against $360$; neither is attained
(Proposition~\ref{prop:values}).

\emph{Roots below $2$.}  The restriction $r_1\ge2$ cannot be dropped once
$u\ge1$, already at $L=2$ (Proposition~\ref{prop:subtwo}).  Keeping the
factors that Lemma~3.1 discards gives the repaired bound \eqref{eq:repaired}
for every $r_1>1$ (Corollary~\ref{cor:repaired}), and a two-sided tail
comparison gives the sharper
$b_{u+1}\ge|f_D|\min(|P(1)|,|P_u(1)|)$ (Corollary~\ref{cor:twosided}), which
improves \eqref{eq:repaired} exactly when $r_1<2<r_u$.  Theorem~\ref{thm:twosidedsharp} and Corollary~\ref{cor:repairedsharp},
repeated roots included (Theorem~\ref{thm:twosidedsharprep}), say exactly
when these bounds are the infimum.  At any distinct roots above $1$
the infimum is the reciprocal of the worst placement's best tail
(Theorem~\ref{thm:value}), but below $2$ the escaping placement need not be
extremal; for $u=1$ the value is still a finite computation
(Proposition~\ref{prop:uoneall}).

\emph{The tail comparison.}  Theorem~2.2 extends to confluent nodes
(Proposition~\ref{prop:confluent}), strictly unless the zero set is
consecutive (Lemma~\ref{lem:confstrict}), which is what the criterion needs
at repeated roots; and its node cutoff $y_i\le1/2$ cannot be raised
(Proposition~\ref{prop:cutoff}).

\paragraph{Organisation.}  Section~\ref{sec:extend} collects the tail
comparisons: the confluent extension, the necessity of the hypotheses, and
the repaired and two-sided bounds.  Section~\ref{sec:criterion} proves the
partial-sum criterion.  Section~\ref{sec:value} computes the infimum when the
criterion fails and proves that the escaping placement is extremal.
Section~\ref{sec:subtwo} decides when the bounds below $2$ are attained.
Section~\ref{sec:scope} lists what remains open and the relation to the
literature.

\section{Tail comparisons and the hypotheses}\label{sec:extend}

\begin{proposition}[Confluent analogue]\label{prop:confluent}
Let $0<y_1<\cdots<y_c\le1/2$ and $e_1,\ldots,e_c\ge1$, put $m=e_1+\cdots+e_c$,
and let $y_{(1)}\le\cdots\le y_{(m)}$ be the multiset with each $y_i$ repeated
$e_i$ times.  Let
\[
  u_d=\sum_{i=1}^c P_i(d)y_i^d,\qquad \deg P_i<e_i,
\]
satisfy $u_0=1$ and $u_z=0$ for a set $Z$ of $m-1$ positive integers with
leading run $f$.  Then
\[
  \Tail(u)\le\prod_{j=1}^{f+1}\frac{y_{(j)}}{1-y_{(j)}}.
\]
\end{proposition}

\begin{proof}
The interpolation system is indexed by the Hermite basis, not by the
ordinary exponentials, and the two must be kept apart.  For
$\varepsilon>0$, $i=1,\ldots,c$ and $k=0,\ldots,e_i-1$ put
\[
  \phi_{i,k,\varepsilon}(d)
  =y_i^d\Bigl(\frac{1-e^{-\varepsilon d}}{\varepsilon}\Bigr)^k .
\]
For fixed $\varepsilon>0$, the binomial expansion gives
\[
  \phi_{i,k,\varepsilon}(d)
  =\frac1{\varepsilon^k}\sum_{j=0}^{k}(-1)^j\binom{k}{j}
    \bigl(y_ie^{-j\varepsilon}\bigr)^d,
\]
so each $\phi_{i,k,\varepsilon}$ is a finite linear combination of the
ordinary exponentials $d\mapsto z^d$ at the distinct nodes
$z_{i,j}=y_ie^{-j\varepsilon}$.  For $0<\varepsilon<\varepsilon_0$ the $m$
nodes $z_{i,j}$ are distinct, lie in $(0,1/2]$, and stay separated between
distinct $i$.  Within one group $i$ the change of basis from
$\{z_{i,j}^{d}\}_{j<e_i}$ to
$\{\phi_{i,k,\varepsilon}\}_{k<e_i}$ is triangular with diagonal entries
$(-1)^k\varepsilon^{-k}$, hence invertible for $\varepsilon>0$, and the same
function space is spanned either way.

For such an $\varepsilon$ let $u^{(\varepsilon)}$ be the unique exponential
sum on the nodes $z_{i,j}$ with $u^{(\varepsilon)}_0=1$ and
$u^{(\varepsilon)}_z=0$ for $z\in Z$; existence and uniqueness are the
Vandermonde argument opening Section~2 of~\cite{coefficient-mass}.  Theorem~2.2 applies with
the same $f$ and gives
\[
  \Tail(u^{(\varepsilon)})
  \le\prod_{j=1}^{f+1}\frac{y_{(j)}(\varepsilon)}{1-y_{(j)}(\varepsilon)},
\]
where $y_{(1)}(\varepsilon)\le\cdots\le y_{(m)}(\varepsilon)$ is the sorted
perturbed multiset.

It remains to pass to the limit.  Write
$a^{(\varepsilon)}=(a^{(\varepsilon)}_{i,k})$ for the coefficients of
$u^{(\varepsilon)}$ in the $\phi$-basis; they solve
$M(\varepsilon)a^{(\varepsilon)}=e_0$, where $M(\varepsilon)$ is the
$m\times m$ matrix with rows indexed by $d\in\{0\}\cup Z$ and columns by
$(i,k)$, with entry $\phi_{i,k,\varepsilon}(d)$, and $e_0$ is the first
standard basis vector.  Each entry
$\phi_{i,k,\varepsilon}(d)=y_i^d\bigl((1-e^{-\varepsilon d})/\varepsilon\bigr)^k$
has a removable singularity at $\varepsilon=0$ and extends analytically there
with limit $d^ky_i^d$, so $M(\varepsilon)\to M(0)$ as $\varepsilon\to0$, where
$M(0)$ is the confluent matrix.  The functions
$d\mapsto d^ky_i^d$ form an extended Chebyshev system on $\R$: a nonzero
$\sum_iP_i(d)y_i^d$ with $\deg P_i<e_i$ has at most $m-1$ real
zeros~\cite[Part~Five, Ch.~1]{polya-szego}.  The
interpolation here is at the $m$ \emph{distinct} points $\{0\}\cup Z$---the
confluence sits in the basis, not in the nodes---so $M(0)$ is invertible by
ordinary unisolvence of a Chebyshev system at distinct
points~\cite[Chapter~I, Theorem~4.1]{karlin-studden}.  Matrix inversion is
continuous on the open set of invertible matrices, so
$a^{(\varepsilon)}\to M(0)^{-1}e_0$, and for every fixed $d$,
\[
  u^{(\varepsilon)}_d
  =\sum_{i,k}a^{(\varepsilon)}_{i,k}\phi_{i,k,\varepsilon}(d)
  \longrightarrow\sum_{i,k}a_{i,k}d^ky_i^d=u_d .
\]
Fatou's lemma for series, applied to the nonnegative terms
$|u^{(\varepsilon)}_d|$, gives
\[
  \Tail(u)=\sum_{d\ge1}|u_d|
  \le\liminf_{\varepsilon\to0}\Tail(u^{(\varepsilon)}).
\]
Finally the order statistics of a converging multiset converge, so the
product $\prod_{j=1}^{f+1}y_{(j)}(\varepsilon)/(1-y_{(j)}(\varepsilon))$
tends to $\prod_{j=1}^{f+1}y_{(j)}/(1-y_{(j)})$.
\end{proof}

\begin{lemma}[Strictness at confluent nodes]\label{lem:confstrict}
In Proposition~\ref{prop:confluent} the inequality is strict unless
$Z=\{1,\ldots,m-1\}$.
\end{lemma}

\begin{proof}
Suppose $Z\ne\{1,\ldots,m-1\}$.  Then $f\le m-2$, the largest element
$z=\max Z$ satisfies $z\ge m$, and $Z'=Z\setminus\{z\}$ has the same leading
run $f$.  Keep the perturbed nodes $z_{i,j}=y_ie^{-j\varepsilon}$ of the
proof of Proposition~\ref{prop:confluent}; the largest of them is
$y_c$ itself ($j=0$), and it is strictly larger than all the others.  Run the
displaced step of the proof of Theorem~2.2 once at these $m$ distinct nodes:
$F^{(\varepsilon)}$ is the sum on the perturbed nodes other than $y_c$ with
$F^{(\varepsilon)}_0=1$ and $F^{(\varepsilon)}|_{Z'}=0$, $G^{(\varepsilon)}$
the sum on all of them with $G^{(\varepsilon)}|_{\{0\}\cup Z'}=0$ and
$G^{(\varepsilon)}_z=1$, and $\mu_\varepsilon=|F^{(\varepsilon)}_z|$.  The
splitting identity and the correction-ratio bound
$\Gamma_G(z)\le2\Gamma_F(z)$ there (all nodes are at most $1/2$) give
\[
  \Tail(u^{(\varepsilon)})
  \le\Tail(F^{(\varepsilon)})-\mu_\varepsilon A_{G^{(\varepsilon)}},
  \qquad
  A_{G^{(\varepsilon)}}=\sum_{1\le d<z}|G^{(\varepsilon)}_d| .
\]
Theorem~2.2 on the $m-1$ nodes of $F^{(\varepsilon)}$ bounds
$\Tail(F^{(\varepsilon)})$ by the product over their $f+1\le m-1$ smallest
members, which are the $f+1$ smallest of all $m$ perturbed nodes because the
one removed is the largest.  So
$\Tail(u^{(\varepsilon)})\le\prod_{j\le f+1}y_{(j)}(\varepsilon)/(1-y_{(j)}(\varepsilon))
-\mu_\varepsilon A_{G^{(\varepsilon)}}$.

Now let $\varepsilon\to0$.  The sums $F^{(\varepsilon)}$ and
$G^{(\varepsilon)}$ are interpolants at the distinct points $\{0\}\cup Z'$
and $\{0\}\cup Z$ in spaces spanned by functions of the form
$\phi_{i,k,\varepsilon}$ (for the group of $y_c$, which has lost its top
member, with $y_c$ replaced by $y_ce^{-\varepsilon}$), so exactly as in the
proof of Proposition~\ref{prop:confluent} they converge pointwise to the
confluent interpolants: $F^{(0)}$ in the Hermite space of the multiset with one
copy of $y_c$ removed, with $F^{(0)}_0=1$ and $F^{(0)}|_{Z'}=0$, and $G^{(0)}$
in the full Hermite space, with $G^{(0)}|_{\{0\}\cup Z'}=0$ and $G^{(0)}_z=1$.
By the P\'olya--Szeg\H o zero bound $F^{(0)}$, a nonzero element of an
$(m-1)$-dimensional Hermite space with the $m-2$ zeros $Z'$, has no other real
zero, so $F^{(0)}_z\ne0$; likewise $G^{(0)}$ vanishes only on
$\{0\}\cup Z'$, and since $|Z'|=m-2<z-1$ some $1\le d<z$ avoids $Z'$, so
$A_{G^{(0)}}>0$.  By Fatou's lemma, as in Proposition~\ref{prop:confluent},
\[
  \Tail(u)\le\liminf_{\varepsilon\to0}\Tail(u^{(\varepsilon)})
  \le\prod_{j=1}^{f+1}\frac{y_{(j)}}{1-y_{(j)}}-|F^{(0)}_z|\,A_{G^{(0)}}
  <\prod_{j=1}^{f+1}\frac{y_{(j)}}{1-y_{(j)}} .
\]
\end{proof}

\begin{lemma}[Coalescence]\label{lem:coalesce}
Let $0<y_1<\cdots<y_c<1$ carry multiplicities $e_i\ge1$, $m=\sum_ie_i$, let
$\mathcal H$ be their Hermite space and $\mathcal H_-$ that of the multiset
with one copy of $y_c$ removed.  For $\varepsilon>0$ let $N_\varepsilon$ be
the multiset of nodes $y_ie^{-j\varepsilon}$, $0\le j<e_i$, and
$N_\varepsilon^-=N_\varepsilon\setminus\{y_c\}$; for small $\varepsilon$ the
members of $N_\varepsilon$ are distinct, lie in $(0,y_c]$, and $y_c$ is the
largest.  Fix a set $Y$ of $m$ (respectively $m-1$) distinct nonnegative
integers and values $b\in\R^Y$, let $v$ be the element of $\mathcal H$
(respectively $\mathcal H_-$) with $v|_Y=b$, and $v^\varepsilon$ the
exponential sum on $N_\varepsilon$ (respectively $N^-_\varepsilon$) with
$v^\varepsilon|_Y=b$.  Then for small $\varepsilon$ there are $C,K$ with
\[
  |v^\varepsilon_d|\le C(1+d)^Ky_c^{\,d}\qquad(d\ge0),
\]
and $v^\varepsilon_d\to v_d$ for every $d$.  Consequently
$\sum_{d\in I}|v^\varepsilon_d|\to\sum_{d\in I}|v_d|$ for every
$I\subseteq\N$; in particular $\Tail(v^\varepsilon)\to\Tail(v)$.
\end{lemma}

\begin{proof}
For $\mathcal H$ this is the limit argument of the proof of
Proposition~\ref{prop:confluent}: in the basis
$\phi_{i,k,\varepsilon}(d)=y_i^d\bigl((1-e^{-\varepsilon d})/\varepsilon\bigr)^k$
the coefficient vectors $a^{(\varepsilon)}$ of $v^\varepsilon$ solve
$M(\varepsilon)a^{(\varepsilon)}=b$ with $M(\varepsilon)\to M(0)$ invertible,
so they converge and are bounded.  For $\mathcal H_-$ the group of $y_c$ has
lost its member $j=0$; the remaining nodes are
$y'e^{-(j-1)\varepsilon}$, $1\le j<e_c$, with $y'=y_ce^{-\varepsilon}$, so
the same construction at the node $y'$ with multiplicity $e_c-1$ gives the
basis $y'^d\bigl((1-e^{-\varepsilon d})/\varepsilon\bigr)^k$, $k<e_c-1$,
whose entries are continuous at $\varepsilon=0$ with limits $d^ky_c^d$, and
the argument is the same.  Since $0\le1-e^{-x}\le x$, every basis function is
bounded by $d^ky_i^d\le d^ky_c^d$ uniformly in $\varepsilon$, which gives the
bound with $K=\max_ie_i-1$ and $C$ the bound on
$\sum|a^{(\varepsilon)}_{i,k}|$.  The dominating sequence is summable because
$y_c<1$, so dominated convergence gives the last claim.
\end{proof}

Every non-strict inequality between tails of interpolants at fixed integer
points, proved for distinct nodes, therefore passes to confluent nodes; the
strict ones need a margin that survives the limit, as in
Lemma~\ref{lem:confstrict}.

\begin{proposition}[The node cutoff $1/2$ is sharp]\label{prop:cutoff}
No constant $a>1/2$ can replace $1/2$ in Theorem~2.2 or in
Proposition~\ref{prop:confluent}: for every $a\in(1/2,1)$ some configuration
with all nodes in $(0,a]$ violates the bound
$\Tail(u)\le\prod_{j=1}^{f+1}y_{(j)}/(1-y_{(j)})$.
\end{proposition}

\begin{proof}
Take a single node $a\in(1/2,1)$ of multiplicity $c$ and
$Z=\{2,3,\ldots,c\}$, so the sum is $u_d=a^d\prod_{k=2}^c(1-d/k)$.  Then
$u_0=1$, the prescribed zeros are exactly $Z$, the leading run is $f=0$, and
\[
  u_1=\frac ac,\qquad
  |u_d|=\frac1c\binom{d-2}{c-1}a^d\quad(d\ge c+1),
\]
so, summing the binomial series,
\[
  \Tail(u)=\frac ac+\frac{a^{c+1}}{c(1-a)^c}
  =\frac ac\Bigl[1+\Bigl(\frac a{1-a}\Bigr)^c\Bigr].
\]
The bound with $f=0$ would be $a/(1-a)$, but
\[
  \frac{\Tail(u)}{a/(1-a)}
  =\frac{1-a}{c}\Bigl[1+\Bigl(\frac a{1-a}\Bigr)^c\Bigr]\longrightarrow\infty
\]
because $a/(1-a)>1$, so $\Tail(u)>a/(1-a)$ for all large $c$.

To transfer the obstruction to distinct nodes, suppose the bound held for
every configuration with nodes in $(0,a]$.  Coalescing the $c$ nodes near $a$
as in the proof of Proposition~\ref{prop:confluent} gives sums
$u^{(\varepsilon)}$ on distinct nodes in $(0,a]$ that satisfy it, and
$u^{(\varepsilon)}_d\to u_d$ for each $d$.  By Fatou's lemma,
$\Tail(u)\le\liminf_{\varepsilon\to0}\Tail(u^{(\varepsilon)})\le a/(1-a)$,
contradicting the previous paragraph for large $c$.
\end{proof}

\begin{proposition}[The root restriction $r_1\ge2$ is necessary for $k\ge2$]
\label{prop:subtwo}
The $k=1$ case of \eqref{eq:order} holds for every root set $r_i>1$ by
Theorem~3.2, and for $L=1$ it is strict: every nonzero multiple
of $x-r$ has $b_1>|f_D|(r-1)$.  For $k\ge2$ the restriction $r_1\ge2$ cannot
be dropped, already at $L=2$, $k=2$, and already when only $r_1<2$: with
$r_1=3/2$ and $r_2=3$ there is a monic multiple $F$ of $(x-r_1)(x-r_2)$ of
degree $4$ with
\[
  b_2=\frac{81}{50}=1.62<2=r_2-1 .
\]
\end{proposition}

\begin{proof}
For $L=1$, $F(r)=0$ gives
$|f_D|r^D=\bigl|\sum_{j<D}f_jr^j\bigr|\le b_1(r^D-1)/(r-1)$, so
$b_1\ge|f_D|(r-1)r^D/(r^D-1)>|f_D|(r-1)$, with no use of $r\ge2$.

For the failure put $P=(x-\tfrac32)(x-3)=x^2+p_1x+p_0$ with
$p_1=-\tfrac92$, $p_0=\tfrac92$, and let $Q=x^2+q_1x+q_0$ be
monic, so that $F=PQ$ has
\[
  f_0=p_0q_0,\quad f_1=p_0q_1+p_1q_0,\quad f_2=p_0+p_1q_1+q_0,\quad
  f_3=p_1+q_1,\quad f_4=1 .
\]
Imposing $f_0=f_1=f_2$ is two linear equations in $(q_0,q_1)$: the first gives
$q_1=(p_0-p_1)q_0/p_0$, and substituting it into the second gives
\[
  q_0=\frac{p_0^2}{\delta},\qquad q_1=\frac{p_0(p_0-p_1)}{\delta},\qquad
  \delta=(p_0-p_1)^2+p_0(p_1-1),
\]
whence the common value is $f_0=f_1=f_2=p_0^3/\delta$.  Here
$\delta=\tfrac{225}4>0$, so
\[
  q_0=\frac9{25},\qquad q_1=\frac{18}{25},\qquad
  f_0=f_1=f_2=\frac{81}{50}=1.62,
\]
while $f_3=p_1+q_1=-\tfrac{189}{50}=-3.78$, that is,
$F=x^4-\tfrac{189}{50}x^3+\tfrac{81}{50}x^2+\tfrac{81}{50}x+\tfrac{81}{50}$.
The four nonleading magnitudes are $|f_3|$ and three copies of
$p_0^3/\delta$, so $b_2=p_0^3/\delta<2=r_2-1$, which is what \eqref{eq:order}
would assert for $u=1$; the second root obeys $r_2\ge2$, so only $r_1$ is out
of range.
\end{proof}

The restriction belongs to the product, though, not to the roots.  The
cutoff $y_i\le1/2$ enters the proof of \eqref{eq:order} in exactly two
places, the factor $2$ in the correction-ratio bound of Theorem~2.2's proof
and the discard of factors at most $1$ in Lemma~3.1, and both matter only for
$u\ge1$; Lemma~2.1, Lemma~2.3 and the case $u=0$ are free of it.  Lemma~3.1
discards, for each exempt position, a factor $y_i/(1-y_i)$ once it is at most
$1$, and that factor exceeds $1$ precisely when $r_i<2$.  Keeping it gives a
bound that holds for every $r_1>1$.

\begin{theorem}[Repaired tail comparison]\label{thm:repaired}
Let $0<y_1<\cdots<y_c<1$, and let $Z$, $f$ and $u$ be as in Theorem~2.2.
Then
\[
  \Tail(u)\le\prod_{i=1}^{f+1}\frac{y_i}{1-y_i}
  \prod_{i=f+2}^{c}\max\Bigl(1,\frac{y_i}{1-y_i}\Bigr).
\]
\end{theorem}

Within the cutoff every factor of the second product is $1$, and this is
Theorem~2.2.

\begin{proof}
Run Theorem~2.2's induction on the number $c-1-f$ of displaced zeros.  The
base case is Lemma~2.1, which needs only $y_i<1$.  In the displaced step, with
$z=\max Z\ge1$ and $F$, $G$, $\mu=|F_z|$ as there, the splitting identity
\[
  \Tail(F)-\Tail(u)=\mu\bigl(2\Gamma_F(z)-\Gamma_G(z)+A_G\bigr),
  \qquad
  \Gamma_V(z)=\frac{\sum_{d\ge z}|V_d|}{|V_z|},
  \quad A_G=\sum_{1\le d<z}|G_d|\ge0,
\]
rests on the zero bound alone, and the first half
$\Gamma_G(z)\le\Gamma_F(z)/(1-y_c)$ of the correction-ratio bound rests on
Lemma~2.3, which needs no upper bound on the nodes, and on a geometric series
in $y_c<1$.  Since $\mu\Gamma_F(z)=\sum_{d\ge z}|F_d|$, together they give
\[
  \Tail(u)\le\Tail(F)+\frac{2y_c-1}{1-y_c}\sum_{d\ge z}|F_d| .
\]
If $y_c\le1/2$ the correction is at most $0$.  Otherwise it is at most
$\frac{2y_c-1}{1-y_c}\Tail(F)$, and $1+\frac{2y_c-1}{1-y_c}=\frac{y_c}{1-y_c}$.
Either way
\[
  \Tail(u)\le\max\Bigl(1,\frac{y_c}{1-y_c}\Bigr)\Tail(F).
\]
The sum $F$ lives on $y_1,\ldots,y_{c-1}$ and its zero set has the same leading
run $f$, so the induction hypothesis bounds $\Tail(F)$ by the product over
those nodes, and the factor for $y_c$ completes the product.
\end{proof}

The deletion step itself can reverse beyond the cutoff: on nodes
$(1/3,1/2,9/10)$ with $Z=\{1,3\}$ it takes $\Tail$ from $2.077$ down to $0.500$.
The factor $\max(1,y_c/(1-y_c))=9$ absorbs the reversal.  The limit argument
of Proposition~\ref{prop:confluent} carries the theorem to confluent nodes
unchanged.  At Proposition~\ref{prop:cutoff}'s witness, one node $a$ of
multiplicity $c$ with $f=0$, it gives $(a/(1-a))^c$ against
$\Tail(u)=(a/c)[1+(a/(1-a))^c]$, a ratio of at most $1/c$.

\begin{corollary}[Repaired order statistics]\label{cor:repaired}
Let $r_1\le\cdots\le r_L$ be real with $r_1>1$, repetitions allowed, and let
$F$ be a nonzero multiple of $\prod_{i=1}^L(x-r_i)$.  Then for every
$0\le u\le L-1$,
\begin{equation}\label{eq:repaired}
  b_{u+1}\ge|f_D|\prod_{i=u+1}^L(r_i-1)\prod_{i=1}^{u}\min(r_i-1,1),
\end{equation}
which is \eqref{eq:order} when $r_u\ge2$.
\end{corollary}

\begin{proof}
Follow Lemma~3.1's proof with Theorem~\ref{thm:repaired} in place of
Theorem~2.2.  The zero set there has leading run $f\ge L-u-1$.  The $L-u$
smallest nodes are the reciprocals of $r_{u+1},\ldots,r_L$ and contribute
$\prod_{i>u}(r_i-1)^{-1}$.  Every other node $1/r_i$, $i\le u$, contributes
either $y_i/(1-y_i)$ or $\max(1,y_i/(1-y_i))$, and both are at most
$1/\min(r_i-1,1)$.  Repeated roots and a general $f_D$ are handled as in
Theorem~3.2's proof.  Its perturbation $r_i+i\varepsilon$ keeps every root
above $1$, and the right side of \eqref{eq:repaired} is continuous in the
roots.
\end{proof}

The repaired bound is itself not the end of the story.  It pays
$\max(1,y_i/(1-y_i))$ for every node outside the leading run, and when the
exempt roots straddle $2$ it pays for the small ones without letting the large
ones pay back.  Keeping the whole product in one piece gives a sharper bound,
and Section~\ref{sec:subtwo} uses it to decide when the infimum is reached.  Write
$q_i=y_i/(1-y_i)$ and $\Pi_j=\prod_{i\le j}q_i$ for nodes
$0<y_1<\cdots<y_c<1$, and for a full certificate $u$ with zero set $Z$ put
$S(u)=\sum_{d>\max Z}|u_d|$, with $\max\varnothing=0$: the mass past the last
zero.

\begin{theorem}[Two-sided tail comparison]\label{thm:twosided}
Let $0<y_1<\cdots<y_c<1$, and let $Z$, $f$ and $u$ be as in Theorem~2.2.  Then
for every $t\ge0$
\[
  \Tail(u)+t\,S(u)\le\max\bigl(\Pi_{f+1},(1+t)\Pi_c\bigr),
\]
with equality exactly when $Z=\{1,\ldots,c-1\}$.  In particular
\[
  \Tail(u)\le\max\Bigl(\prod_{i=1}^{f+1}\frac{y_i}{1-y_i},\
                       \prod_{i=1}^{c}\frac{y_i}{1-y_i}\Bigr),
\]
strictly unless $Z$ is consecutive.
\end{theorem}

Since $\Pi_c\le\Pi_{f+1}\prod_{i>f+1}\max(1,q_i)$, this sharpens
Theorem~\ref{thm:repaired}; within the cutoff every $q_i\le1$, the maximum is
$\Pi_{f+1}$, and it is Theorem~2.2.

\begin{proof}
Induct on the number $c-1-f$ of displaced zeros.  If $Z=\{1,\ldots,c-1\}$
(including $c=1$, $Z=\varnothing$), then $f=c-1$, $u_d=0$ for $1\le d<c$, and
Lemma~2.1 gives $S(u)=\Tail(u)=\Pi_c$, so both sides equal $(1+t)\Pi_c$.

Otherwise let $z=\max Z$, $Z'=Z\setminus\{z\}$, and let $F$, $G$ and
$\mu=|F_z|>0$ be as in the displaced step of Theorem~2.2's proof: $F$ is the
full certificate on $y_1,\ldots,y_{c-1}$ with zero set $Z'$, which has the same
leading run $f\le c-2$.  The sign comparison there uses only the zero bound and
gives $|u_d|=\mu|G_d|-|F_d|$ for $d>z$, so, with $|G_z|=1$ and
$\mu\Gamma_F(z)=\sum_{d\ge z}|F_d|$,
\[
  S(u)=\mu\bigl(\Gamma_G(z)-1\bigr)-\bigl(\mu\Gamma_F(z)-\mu\bigr)
      =\mu\bigl(\Gamma_G(z)-\Gamma_F(z)\bigr).
\]
Adding $t$ times this to the splitting identity of Theorem~2.2's proof,
\[
  \Tail(u)+t\,S(u)=\Tail(F)+\mu\bigl((1+t)\Gamma_G(z)-(2+t)\Gamma_F(z)
  -\Sigma_G\bigr).
\]
The first half $\Gamma_G(z)\le\Gamma_F(z)/(1-y_c)=(1+q_c)\Gamma_F(z)$ of the
correction-ratio bound needs no cutoff, and $\Sigma_G>0$ by the strictness
paragraph of that proof, which needs none either.  With $t'=q_c(1+t)-1$,
\[
  \Tail(u)+t\,S(u)<\Tail(F)+t'\sum_{d\ge z}|F_d| .
\]
If $t'\ge0$, then $z>\max Z'$ gives $\sum_{d\ge z}|F_d|\le S(F)$, and the
induction hypothesis for $F$ at the parameter $t'$ bounds the right side by
$\max(\Pi_{f+1},(1+t')\Pi_{c-1})$; and $(1+t')\Pi_{c-1}=(1+t)q_c\Pi_{c-1}
=(1+t)\Pi_c$.  If $t'<0$, then $q_c<1$, so every $q_i<1$ and
$\Pi_{c-1}\le\Pi_{f+1}$; the right side is at most $\Tail(F)$, which the
induction hypothesis at $t=0$ bounds by $\max(\Pi_{f+1},\Pi_{c-1})=\Pi_{f+1}$.
Either way the bound holds, strictly.
\end{proof}

For a multiset $y_{(1)}\le\cdots\le y_{(m)}$ in $(0,1)$ put
$q_{(j)}=y_{(j)}/(1-y_{(j)})$ and $\Pi_j=\prod_{i\le j}q_{(i)}$, and for a
certificate $u$ with zero set $Z$ put $S(u)=\sum_{d>\max Z}|u_d|$.

\begin{proposition}[Two-sided comparison at confluent nodes]
\label{prop:twosidedconf}
Let $0<y_1<\cdots<y_c<1$ with multiplicities $e_i$, $m=\sum_ie_i$, and let
$u$ be the element of their Hermite space with $u_0=1$ and $u|_Z=0$ for a
set $Z$ of $m-1$ positive integers with leading run $f$.  Then for every
$t\ge0$
\[
  \Tail(u)+t\,S(u)\le\max\bigl(\Pi_{f+1},(1+t)\Pi_m\bigr),
\]
with equality exactly when $Z=\{1,\ldots,m-1\}$.
\end{proposition}

\begin{proof}
If $Z=\{1,\ldots,m-1\}$, Lemma~2.1 with repetitions (the paragraph after
Lemma~\ref{lem:improve}) gives $u_d=0$ for $1\le d<m$, one sign past $m-1$,
and $S(u)=\Tail(u)=\Pi_m=\Pi_{f+1}$; both sides are $(1+t)\Pi_m$.

Otherwise $z=\max Z\ge m$ and $Z'=Z\setminus\{z\}$ has the same leading run
$f\le m-2$.  Take $N_\varepsilon$ as in Lemma~\ref{lem:coalesce}, let
$u^{(\varepsilon)}$ be the sum on $N_\varepsilon$ with the conditions of
$u$, and run the displaced step of the proof of Theorem~\ref{thm:twosided}
at these $m$ distinct nodes, deleting $z$ together with the largest node
$y_c$: $F^{(\varepsilon)}$ is the sum on $N^-_\varepsilon$ with
$F^{(\varepsilon)}_0=1$ and $F^{(\varepsilon)}|_{Z'}=0$, $G^{(\varepsilon)}$
the sum on $N_\varepsilon$ with $G^{(\varepsilon)}|_{\{0\}\cup Z'}=0$ and
$G^{(\varepsilon)}_z=1$, $\mu_\varepsilon=|F^{(\varepsilon)}_z|$ and
$\Sigma_{G}=\sum_{1\le d<z}|G_d|$.  The identity of that proof and the first
half of the correction-ratio bound, which needs no cutoff, give with
$t'=q_c(1+t)-1$, $q_c=y_c/(1-y_c)$,
\[
  \Tail(u^{(\varepsilon)})+t\,S(u^{(\varepsilon)})
  \le\Tail(F^{(\varepsilon)})+t'\sum_{d\ge z}|F^{(\varepsilon)}_d|
     -\mu_\varepsilon\Sigma_{G^{(\varepsilon)}} .
\]
The nodes of $F^{(\varepsilon)}$ are the $m-1$ smallest members of
$N_\varepsilon$, so, writing $\Pi_j(\varepsilon)$ for the products over the
sorted $N_\varepsilon$, Theorem~\ref{thm:twosided} for $F^{(\varepsilon)}$
bounds the first two terms by $\max(\Pi_{f+1}(\varepsilon),
(1+t')\Pi_{m-1}(\varepsilon))=\max(\Pi_{f+1}(\varepsilon),(1+t)\Pi_m(\varepsilon))$
if $t'\ge0$, and by $\Pi_{f+1}(\varepsilon)$ if $t'<0$, exactly as in that
proof.  Hence
\[
  \Tail(u^{(\varepsilon)})+t\,S(u^{(\varepsilon)})
  \le\max\bigl(\Pi_{f+1}(\varepsilon),(1+t)\Pi_m(\varepsilon)\bigr)
     -\mu_\varepsilon\Sigma_{G^{(\varepsilon)}} .
\]
Let $\varepsilon\to0$.  By Lemma~\ref{lem:coalesce} (with $I=\N_{>0}$ and
$I=(z,\infty)$) the left side tends to $\Tail(u)+tS(u)$;
$F^{(\varepsilon)}\to F^{(0)}$, the element of the Hermite space of the
multiset with one copy of $y_c$ removed with $F^{(0)}_0=1$,
$F^{(0)}|_{Z'}=0$; and $G^{(\varepsilon)}\to G^{(0)}$ pointwise.  As in the
proof of Lemma~\ref{lem:confstrict}, the zero bound gives $F^{(0)}_z\ne0$,
and $G^{(0)}$ vanishes only on $\{0\}\cup Z'$, while $|Z'|=m-2<z-1$ leaves
some $1\le d<z$ outside $Z'$; so the finite sums
$\mu_\varepsilon\Sigma_{G^{(\varepsilon)}}$ tend to
$|F^{(0)}_z|\Sigma_{G^{(0)}}>0$.  The products converge, and the inequality
is strict in the limit.
\end{proof}

With $t=0$ and all nodes at most $1/2$, where $\Pi_{f+1}\ge\Pi_m$, this is
Proposition~\ref{prop:confluent} with Lemma~\ref{lem:confstrict}.  Read on
root sets as in the proof of Corollary~\ref{cor:twosided}, it gives
\eqref{eq:twosided} at repeated roots without perturbing the roots.  Read on root sets, the maximum of
two products becomes the minimum of two root products.

\begin{corollary}[Two-sided order statistics]\label{cor:twosided}
Let $r_1\le\cdots\le r_L$ be real with $r_1>1$, repetitions allowed, put
$P=\prod_{i=1}^L(x-r_i)$ and $P_u=\prod_{i>u}(x-r_i)$, and let $F$ be a nonzero
multiple of $P$.  Then for every $0\le u\le L-1$
\begin{equation}\label{eq:twosided}
  b_{u+1}\ge|f_D|\min\bigl(|P(1)|,|P_u(1)|\bigr)
  =|f_D|\prod_{i=u+1}^L(r_i-1)\cdot\min\Bigl(1,\prod_{i=1}^u(r_i-1)\Bigr).
\end{equation}
This implies \eqref{eq:repaired}, and exceeds it exactly when
$r_1<2<r_u$.
\end{corollary}

\begin{proof}
For distinct roots and $|f_D|=1$ follow Lemma~3.1's proof with
Theorem~\ref{thm:twosided} in place of Theorem~2.2.  In increasing order the
nodes are $1/r_L<\cdots<1/r_1$, so
$\Pi_j=\prod_{i>L-j}(r_i-1)^{-1}$, and the zero set there has leading run
$f\ge m-1$, $m=L-u$.  The increments $\log\Pi_j-\log\Pi_{j-1}=
-\log(r_{L-j+1}-1)$ increase with $j$, so $j\mapsto\log\Pi_j$ is convex and,
as $m\le f+1\le L$, $\Pi_{f+1}\le\max(\Pi_m,\Pi_L)$.  Hence
$\Tail\le\max(\Pi_m,\Pi_L)=\max(1/|P_u(1)|,1/|P(1)|)$, which is
\eqref{eq:twosided}.  A general $f_D$ and repeated roots are handled as in
Theorem~3.2's proof; the right side is continuous in the roots.

Put $s_i=r_i-1$ and let $p$ be the number of $i\le u$ with $s_i<1$.  The
factor in \eqref{eq:repaired} is
$\prod_{i\le u}\min(s_i,1)=\prod_{i\le p}s_i$, that in \eqref{eq:twosided} is
$\min(1,\prod_{i\le u}s_i)$.  They agree if $p=0$ and if all $s_i\le1$.
Otherwise $r_1<2<r_u$, the first is below $1$, and it is below
$\prod_{i\le u}s_i$ by the factor $\prod_{p<i\le u}s_i\ge s_u>1$.
\end{proof}

Section~\ref{sec:subtwo} decides when \eqref{eq:repaired} and
\eqref{eq:twosided} are attained.

\section{The partial-sum criterion}\label{sec:criterion}

This section proves that \eqref{eq:order} is the infimum exactly when the
partial-sum condition \eqref{eq:partial} holds.  The condition comes out of
an explicit construction, which we first carry out at $(2,3,5)$.

At roots $(2,3)$ the certificate paper already settles the question: it
exhibits monic multiples
$F_n=x^{n+1}-a_nx^n+t_n\sum_{j=0}^{n-1}x^j$ of $(x-2)(x-3)$ with
$b_2(F_n)=t_n$ decreasing to $2$~\cite[Proposition~4.5]{coefficient-mass}, so
the $u=1$ bound is an infimum there, approached but never attained.

The same construction reaches $L=3$.  Write the multiple from the top: one
leading $1$, a short free prefix, a long run of a single value $t$, and $u$
exempt coefficients at the bottom.  The exempt coefficients grow like
$r_u^D$ and so survive division by $r^{D-1}$ only at the $u$ smallest roots,
which they absorb; the run must therefore satisfy the remaining $L-u$ roots
alone, and its value is what \eqref{eq:order} predicts.  The certificate
paper proves the two limits below~\cite[Section~4]{coefficient-mass}; what
is added here is that the ordering holds exactly at every degree, with
$t_D>a_D+1>8$ throughout, and the Descartes step that the general
construction later reuses.

\begin{proposition}[The bounds at roots $(2,3,5)$ are infima]
\label{prop:sharp235}
For $D\ge5$ let $(a_D,t_D,c_D)$ be the solution of the linear system
$F_D(2)=F_D(3)=F_D(5)=0$ in
\[
  F_D(x)=x^D-a_Dx^{D-1}+t_D\sum_{j=1}^{D-2}x^j+c_D ,
\]
and for $D\ge6$ let $(s_D,e_D,e_D')$ be the solution of
$G_D(2)=G_D(3)=G_D(5)=0$ in
\[
  G_D(x)=x^D-s_D\sum_{j=2}^{D-1}x^j+e_Dx+e_D' .
\]
Both systems are nonsingular, $F_D$ and $G_D$ are monic multiples of
$(x-2)(x-3)(x-5)$ with $b_2(F_D)=t_D$ and $b_3(G_D)=s_D$, and $t_D\to8$,
$s_D\to4$.  Hence
\begin{equation}\label{eq:sharp235}
  \inf_Q b_2\bigl(QP\bigr)=8=(r_2-1)(r_3-1),
  \qquad
  \inf_Q b_3\bigl(QP\bigr)=4=r_3-1,
\end{equation}
over monic $Q\in\R[x]$, with $P(x)=(x-2)(x-3)(x-5)$: the values
\eqref{eq:order} asserts for $u=1$ and $u=2$.
\end{proposition}

\begin{proof}
Both systems are nonsingular for every $D$ in range.  A kernel vector of the
first is an $H=-ax^{D-1}+t\sum_{j=1}^{D-2}x^j+c$ vanishing at $2,3,5$, and
\[
  (x-1)H=-ax^D+(a+t)x^{D-1}+(c-t)x-c
\]
has four terms, hence at most three positive roots by Descartes' rule; but
$1,2,3,5$ are four, so $H\equiv0$ and $a=t=c=0$.  The second is the same with
$(x-1)H=-sx^D+(s+e)x^2+(e'-e)x-e'$.

Divide $F_D(r)=0$ by $r^{D-1}$ and sum the geometric progression:
\[
  (r-a_D)+\frac{t_D}{r-1}
  +r^{1-D}\Bigl(c_D-\frac{r\,t_D}{r-1}\Bigr)=0,
  \qquad r\in\{2,3,5\}.
\]
Put $\gamma_D=2^{1-D}c_D$, so that $r^{1-D}c_D=\gamma_D(2/r)^{D-1}$.  In the
unknowns $(a_D,t_D,\gamma_D)$ the three equations read
\[
  \begin{aligned}
    -a+t\bigl(1-2^{2-D}\bigr)+\gamma&=-2,\\
    -a+t\bigl(\tfrac12-\tfrac323^{1-D}\bigr)+\gamma(2/3)^{D-1}&=-3,\\
    -a+t\bigl(\tfrac14-\tfrac545^{1-D}\bigr)+\gamma(2/5)^{D-1}&=-5 .
  \end{aligned}
\]
Every $D$-dependent entry is $O((2/3)^{D})$, so the coefficient matrix tends
to
\[
  M=\begin{pmatrix}-1&1&1\\-1&\tfrac12&0\\-1&\tfrac14&0\end{pmatrix},
  \qquad \det M=\tfrac14\ne0 .
\]
Its inverse therefore converges to
$M^{-1}$, so $(a_D,t_D,\gamma_D)$ converges to the solution of $M$ against
$(-2,-3,-5)$.  The last two rows of that limit give
$-a+\tfrac12t=-3$ and $-a+\tfrac14t=-5$, whence $t=8$ and $a=7$; the first
row then gives $\gamma=-3$, so $c_D\sim-3\cdot2^{D-1}$.

The ordering is exact for every $D$, not only in the limit.  $F_D$ has the
three positive roots $2,3,5$, and its coefficient sequence
$(1,-a_D,t_D,\ldots,t_D,c_D)$ offers three sign changes only if $a_D>0$,
$t_D>0$ and $c_D<0$, a vanishing entry leaving at most two; so Descartes
forces all three signs.  Multiplying the displayed equation at $r$ by $r-1$
gives $(r-1)(r-a_D)+t_D=\varepsilon_r$ with
$\varepsilon_r=r^{1-D}\bigl(r\,t_D+(r-1)|c_D|\bigr)>0$, and the rows $r=3,5$
solve to
\[
  a_D=7+\tfrac12(\varepsilon_3-\varepsilon_5),\qquad
  t_D=8+2\varepsilon_3-\varepsilon_5,\qquad
  t_D-a_D=1+\tfrac32\varepsilon_3-\tfrac12\varepsilon_5 ,
\]
where
$\varepsilon_5/\varepsilon_3=(3/5)^{D-1}(5t_D+4|c_D|)/(3t_D+2|c_D|)
<2(3/5)^{D-1}<1$.  Hence $t_D>a_D+1>8$: both limits are approached from
above.  The nonleading magnitudes are $a_D$ once, $t_D$ exactly $D-2\ge3$
times and $|c_D|$ once, and the run already exceeds $a_D$, so at most one
entry exceeds $t_D$ and $b_2(F_D)=t_D$.  Theorem~3.2 with $u=1$
gives $b_2\ge8$ for every monic multiple, and $t_D\to8$.

For $G_D$ the two exempt coefficients grow like $3^D$ and must absorb both
$2$ and $3$, so the rescaling is the one that keeps the equation at $r=2$
finite: put
\[
  E_D=3^{2-D}e_D,
  \qquad
  \sigma_D=2^{1-D}\bigl(2e_D+e_D'\bigr).
\]
Dividing $G_D(r)=0$ by $r^{D-1}$ and substituting, the exempt terms
contribute $\sigma_D$ at $r=2$ by definition, $E_D/3+\sigma_D(2/3)^{D-1}$ at
$r=3$, and $O((3/5)^D)$ at $r=5$.  In the unknowns $(s_D,E_D,\sigma_D)$ the
system therefore reads
\[
  \begin{aligned}
    -2s\bigl(1-2^{2-D}\bigr)+\sigma&=-2,\\
    -\tfrac32s\bigl(1-3^{2-D}\bigr)+\tfrac13E+\sigma(2/3)^{D-1}&=-3,\\
    -\tfrac54s\bigl(1-5^{2-D}\bigr)+O\bigl((3/5)^D\bigr)&=-5,
  \end{aligned}
\]
whose coefficient matrix tends to
\[
  N=\begin{pmatrix}-2&0&1\\-\tfrac32&\tfrac13&0\\-\tfrac54&0&0\end{pmatrix},
  \qquad \det N=\tfrac5{12}\ne0 .
\]
Every $D$-dependent entry is $O((2/3)^{D})$, so as before the solution
converges to that of $N$ against $(-2,-3,-5)$, namely
$(s,E,\sigma)=(4,9,6)$; the last row alone forces $s=4$.  The nonleading
magnitudes of $G_D$ are $s_D$ exactly $D-2$ times and $|e_D|,|e_D'|$ once
each, and the same sign count applies: $(1,-s_D,\ldots,-s_D,e_D,e_D')$ offers
three changes only if $s_D>0$, so $s_D>0$ for every $D$.  The run then
occupies $D-2\ge3$ positions, and whatever the two exempt magnitudes do the
third largest is $s_D$, so $b_3(G_D)=s_D$.
Theorem~3.2 with $u=2$ gives $b_3\ge4$, and
$s_D\to4$, which completes \eqref{eq:sharp235}.
\end{proof}

Neither family was guessed.  Write $P_u(x)=\prod_{i>u}(x-r_i)$ for the
surviving roots, $m=L-u$, and let $\widetilde P_u(y)=y^mP_u(1/y)$ be its
reverse.  In the reversed variables the limit construction is
\begin{equation}\label{eq:wgen}
  W(y)=\sum_{d\ge0}w_dy^d=\frac{\widetilde P_u(y)}{1-y},
\end{equation}
which vanishes at $y=1/r_i$ for every surviving root because
$\widetilde P_u$ does.  Its coefficients are therefore the partial sums
$S_j=c_m+c_{m-1}+\cdots+c_{m-j}$ of the coefficients $c_k$ of $P_u$ read down
from the leading $c_m=1$, and they stabilize at
$S_m=P_u(1)=(-1)^m\prod_{i>u}(r_i-1)$ once $j\ge m$.  So the free prefix is
$S_1,\ldots,S_{m-1}$ and the run value is $|S_m|$, which is exactly
\eqref{eq:order}: the construction reproduces the bound because the sum of the
coefficients of $P_u$ is $P_u(1)$.

This makes the requirement explicit.  At $u=0$ both the criterion and this
series are already in Dubickas and Jankauskas: their Corollary~7 gives
$H(P)=|P(1)|$ under the same partial-sum condition, together with the one-sign
proviso, and exhibits $\widetilde P/(1-y)$ as the extremal
series~\cite[Corollary~7]{dj}.  What is added here is the passage to $u\ge1$,
where the $u$ exempt coefficients absorb the $u$ smallest roots.  The run is
the $(u+1)$-st largest magnitude only if no partial sum exceeds the total,
\begin{equation}\label{eq:partial}
  |S_j|\le|S_m|\qquad(1\le j<m),
\end{equation}
and \eqref{eq:partial} can fail.

At $(2,3,5,7)$ with $u=1$ the coefficients
of $P_1$ are $1,-15,71,-105$, with partial sums $1,-14,57,-48$; since
$57>48$, \eqref{eq:partial} fails.  The same happens at $(2,3,4,5)$, sums
$1,-11,36,-24$, and at $(2,3,5,6)$, sums $1,-13,50,-40$.  At $(3,4,5,6)$, by
contrast, the sums are $1,-14,60,-60$, so \eqref{eq:partial} holds with
equality.  It can fail already at $L=3$: at $(2,3,\tfrac72)$ with $u=1$ the
partial sums are $1,-\tfrac{11}2,5$, so by Theorem~\ref{thm:converse} below
the infimum of $b_2$ exceeds $5$.

\begin{proposition}[Partial-sum criterion]\label{prop:partial}
Let $2\le r_1<\cdots<r_L$ be distinct and $0\le u\le L-1$.  If
\eqref{eq:partial} holds, then
$\inf_Qb_{u+1}\bigl(Q\prod_{i=1}^L(x-r_i)\bigr)=\prod_{i=u+1}^L(r_i-1)$ over
monic $Q\in\R[x]$.
\end{proposition}

\begin{proof}
Theorem~3.2 gives $\ge$.  For the other direction put $m=L-u$ and write
$F_D=x^D+\sum_{k=1}^{m-1}a_kx^{D-k}+t\sum_{j=u}^{D-m}x^j+C$ with $\deg C<u$,
the $L$ unknowns fixed by $F_D(r_i)=0$.  The Descartes step of
Proposition~\ref{prop:sharp235} carries over, $(x-1)$ times a kernel element
having $L+1$ terms against the $L+1$ positive roots $1,r_1,\ldots,r_L$.
Divide by $r_i^D$, sum the run, and trade the exempt coefficients for their
values $\sigma_i=r_i^{-D}C(r_i)$, a Vandermonde change of variables.  The
limit system is block triangular, and its surviving rows make the monic
$(x-1)\widetilde A+t$ of degree $m$, where
$\widetilde A=x^{m-1}+a_1x^{m-2}+\cdots+a_{m-1}$, vanish at the $m$ surviving
roots.  So it is $P_u$, and
\[
  t=P_u(1),\qquad
  \widetilde A=\frac{P_u-P_u(1)}{x-1},\qquad
  \sigma_i=-\frac{r_i^{1-m}P_u(r_i)}{r_i-1}\ne0 ,
\]
the coefficients of $\widetilde A$ being the prefix $S_1,\ldots,S_{m-1}$ of
\eqref{eq:wgen}.  The perturbation is $O((r_u/r_{u+1})^D)$ and every exempt
coefficient grows like $r_u^D$, so for large $D$ those are the $u$ largest and
$b_{u+1}\to\max_j|S_j|$, which is $|S_m|$ exactly when \eqref{eq:partial}
holds.  Distinctness enters three times---the Descartes step, the Vandermonde
change, and the surviving block---and $r_u<r_{u+1}$ is what kills the cross
terms and makes the exempt coefficients diverge.
\end{proof}

The conclusion does not need distinct roots (Corollary~\ref{cor:charrep}),
and at repeated roots the approach can be slower.  At $(r,r)$ the
confluent member $F_D=x^D-t_D\sum_{j=1}^{D-1}x^j+c_D$ fixed by
$F_D(r)=F_D'(r)=0$ has $b_2(F_D)=t_D$, the run occupying $D-1\ge2$ positions,
and at $r=3$
\[
  t_D=\frac{4D\cdot3^{D-1}}{3^{D-1}(2D-3)+1}\longrightarrow2=r_2-1,
\]
so \eqref{eq:order} is an infimum there too, approached at rate $1/D$ rather
than geometrically.

\subsection*{Necessity}

The criterion is also necessary, repeated roots included.  The argument runs
on the dual side of Lemma~3.1; we set it up for distinct roots, and the proof
of Theorem~\ref{thm:converse} carries it to repeated roots.  Put $y_i=1/r_i$, so that
$1/2\ge y_1>\cdots>y_L>0$, and call an exponential sum
$v_d=\sum_i\lambda_iy_i^d$ on some of these nodes a \emph{certificate for}
$X\subset\N_{>0}$ if $v_0=1$ and $v_w=0$ for $w\in X$.  If $F$ is a monic
multiple of $\prod_{i=1}^L(x-r_i)$ of degree $D$ and $X$ is the set of
distances $D-j$ of $u$ positions $j$ carrying its $u$ largest nonleading
magnitudes, the identity in the proof of Lemma~3.1 reads
$0=1+\sum_{j<D}f_jv_{D-j}$, the terms with $D-j\in X$ vanish, and every other
$|f_j|$ is at most $b_{u+1}(F)$, so
\begin{equation}\label{eq:weak}
  1\le b_{u+1}(F)\,\Tail(v)
\end{equation}
for every certificate $v$ for $X$.  Write
$B=\prod_{i=u+1}^L(r_i-1)^{-1}$.  It therefore suffices to find a
$\delta>0$ such that every $X$ of size $u$ has a certificate with
$\Tail(v)\le B-\delta$.  Three ingredients supply the pieces: the partial
sums alternate in sign (Lemma~\ref{lem:signs}), a failing partial sum
improves the consecutive certificate (Lemma~\ref{lem:improve}), and zeros far
from the origin cost only the largest nodes
(Proposition~\ref{prop:farrep}).

\begin{lemma}[Signs of the partial sums]\label{lem:signs}
Let $r_1,\ldots,r_m>1$, let $S_0,\ldots,S_m$ be the partial sums of the
coefficients of $\prod_{i=1}^m(x-r_i)$ read down from the leading one, and put
$s_i=r_i-1$.  Then
\[
  (-1)^jS_j=\sum_{k=0}^{j}e_k(s)\binom{m-1-k}{j-k}>0\quad(0\le j\le m-1),
  \qquad (-1)^mS_m=e_m(s),
\]
where $e_k$ is the elementary symmetric polynomial.
\end{lemma}

\begin{proof}
By \eqref{eq:wgen}, $S_j$ is the coefficient of $y^j$ in
$\prod_i(1-r_iy)/(1-y)$, so $(-1)^jS_j$ is the coefficient of $y^j$ in
$\prod_i(1+r_iy)/(1+y)$.  Writing $1+r_iy=(1+y)+s_iy$ and grouping the
expansion by the number of factors $s_iy$ chosen,
\[
  \frac{\prod_i(1+r_iy)}{1+y}
  =\sum_{k=0}^{m-1}e_k(s)\,y^k(1+y)^{m-1-k}+e_m(s)\frac{y^m}{1+y}.
\]
The terms with $k<m$ are polynomials of degree $m-1$ and contribute
$e_k(s)\binom{m-1-k}{j-k}$ to $y^j$; the last contributes nothing below $y^m$
and $e_m(s)$ at $y^m$.  The $k=0$ term alone gives $\binom{m-1}{j}\ge1$.
\end{proof}

In particular \eqref{eq:partial} reads
$\sum_{k\le j}e_k(s)\binom{m-1-k}{j-k}\le e_m(s)$ for $1\le j<m$, with
$s_i=r_i-1$ over the surviving roots; at $(3,5,7)$, $s=(2,4,6)$, the two sums
are $14$ and $57$ against $e_3(s)=48$.

\begin{lemma}[Improving the consecutive certificate]\label{lem:improve}
Let $\rho_1,\ldots,\rho_n>1$ be distinct, $Q=\prod_{i=1}^n(x-\rho_i)$ with
partial sums $S_j(Q)$, and $v^*$ the certificate on the nodes $1/\rho_i$ with
zero set $\{1,\ldots,n-1\}$, so that $\Tail(v^*)=1/|Q(1)|$ by Lemma~2.1.  If
$X_0\subseteq\{1,\ldots,n-1\}$ and some $j\in\{1,\ldots,n-1\}\setminus X_0$
has $|S_j(Q)|>|Q(1)|$, then some certificate for $X_0$ on the same nodes has
tail less than $1/|Q(1)|$.
\end{lemma}

\begin{proof}
Let $\delta$ be the exponential sum on the nodes with $\delta_0=0$ and
$\delta_i=[i=j]$ for $1\le i\le n-1$; it exists by the Vandermonde argument.
With $\Pi(x)=\prod_i(1-x/\rho_i)=\sum_s\pi_sx^s$, its generating function is
$N(x)/\Pi(x)$ with $\deg N\le n-1$, and $N$ is the truncation of
$\Pi\sum_d\delta_dx^d$ below degree $n$, namely
$N(x)=\sum_{t=j}^{n-1}\pi_{t-j}x^t$.  Since $\Pi$ is $Q$ times the constant
$\prod_i(-1/\rho_i)$, evaluating at $x=1$ gives
\[
  \sum_{d\ge n}\delta_d=\frac{N(1)}{\Pi(1)}-1
  =-\frac{\pi_{n-j}+\cdots+\pi_n}{\Pi(1)}
  =-\frac{S_j(Q)}{Q(1)} .
\]
By Lemma~2.1, $v^*_d$ has the constant sign $\sigma=(-1)^{n+1}$ for $d\ge n$
and $|v^*_d|\ge\bigl(\prod_iy_i\bigr)y_{\max}^{\,d-n}$, while
$|\delta_d|\le Cy_{\max}^{\,d}$; so for $0<\tau$ small enough
$v=v^*+t\delta$ with $t=\tau\sigma\sgn\bigl(S_j(Q)/Q(1)\bigr)$ keeps the sign
$\sigma$ at every $d\ge n$.  It is a certificate for $X_0$, since it vanishes
on $\{1,\ldots,n-1\}\setminus\{j\}$, and
\[
  \Tail(v)=\tau+\sigma\sum_{d\ge n}\bigl(v^*_d+t\delta_d\bigr)
  =\Tail(v^*)-\tau\Bigl(\Bigl|\frac{S_j(Q)}{Q(1)}\Bigr|-1\Bigr)
  <\Tail(v^*).
\]
\end{proof}

\paragraph{Lemma~\ref{lem:improve} with repetitions.}  The lemma holds with
repetitions allowed among the $\rho_i$, the nodes being replaced by the
Hermite space of $Q$.  The proof of Lemma~2.1 is a generating-function
computation that uses only the denominator $\prod_i(1-y_ix)$, so with
repetitions it still gives
$v^*_d=(-1)^{n+1}(\prod_iy_i)h_{d-n}(y)$ for $d\ge n$ and
$\Tail(v^*)=1/|Q(1)|$; the sum $\delta$ is given explicitly by its generating
function $N(x)/\Pi(x)$, and the evaluation at $x=1$ is unchanged.  The only
change is in the domination: if $y_{\max}$ has multiplicity $\mu$, then
$h_{d-n}(y)\ge\binom{d-n+\mu-1}{\mu-1}y_{\max}^{\,d-n}$ while
$|\delta_d|\le C(1+d)^{\mu-1}y_{\max}^{\,d}$, so $|\delta_d/v^*_d|$ is still
bounded for $d\ge n$, and a small $\tau$ keeps the sign of $v^*$.

Zeros far from the origin cost only the largest nodes because the leading
functions of the basis cannot produce them, and that is a Descartes rule at
infinity.  Let $\mathcal C$ be the
set of functions $x\mapsto R(x)w^x$ with $R\in\R(x)$ nonzero and $w>0$, each
considered on a half-line beyond the real poles of $R$.  It is a group under
multiplication; each member has finitely many zeros, so it is eventually of
one sign; the quotient of two members,
$(R_1/R_2)(w_1/w_2)^x$, has a limit in $[-\infty,\infty]$ as $x\to\infty$;
and the derivative of a nonconstant member lies in $\mathcal C$, because
$(Rw^x)'=(R'+R\log w)w^x$ and, writing $R=p/q$, the numerator
$p'q-pq'+pq\log w$ of $R'+R\log w$ has a term of degree $\deg p+\deg q$ that
the other two cannot cancel when $w\ne1$, while $R'\ne0$ for nonconstant $R$
when $w=1$.

\begin{lemma}[Descartes at infinity]\label{lem:infinity}
Let $\varphi_1,\ldots,\varphi_M\in\mathcal C$ satisfy
$\varphi_{j+1}(x)/\varphi_j(x)\to0$ as $x\to\infty$.  Let
$c^{(\nu)}\in\R^M$ converge to $c\ne0$, let $j_0$ be the least index with
$c_{j_0}\ne0$, and let $n_\nu\to\infty$.  Then for all large $\nu$ the
function $f_\nu=\sum_jc^{(\nu)}_j\varphi_j$ has at most $j_0-1$ zeros in
$[n_\nu,\infty)$.
\end{lemma}

\begin{proof}
Induct on $j_0$, for all such families at once.  Every ratio
$\varphi_j/\varphi_1$, $j\ge2$, is a product of consecutive ratios and tends
to $0$.  If $j_0=1$ then on $[n_\nu,\infty)$
\[
  \Bigl|\frac{f_\nu}{\varphi_1}-c^{(\nu)}_1\Bigr|
  \le\max_j|c^{(\nu)}_j|\sum_{j\ge2}\sup_{x\ge n_\nu}
  \Bigl|\frac{\varphi_j(x)}{\varphi_1(x)}\Bigr|\longrightarrow0,
\]
while $c^{(\nu)}_1\to c_1\ne0$, so $f_\nu$ has no zero there for large $\nu$.

If $j_0\ge2$, then $c_1=0$.  Put $\chi_j=(\varphi_j/\varphi_1)'$ for
$2\le j\le M$.  Each $\varphi_j/\varphi_1$ tends to $0$ and is not constant,
so $\chi_j\in\mathcal C$.  For $2\le j<M$ the functions
$\varphi_{j+1}/\varphi_1$ and $\varphi_j/\varphi_1$ tend to $0$, their
quotient $\varphi_{j+1}/\varphi_j$ tends to $0$, and the quotient
$\chi_{j+1}/\chi_j$ of their derivatives lies in $\mathcal C$ and so has a
limit $\ell\in[-\infty,\infty]$, while $\chi_j$ is eventually nonzero.
L'H\^opital's rule then gives $\varphi_{j+1}/\varphi_j\to\ell$, so $\ell=0$:
the family $\chi_2,\ldots,\chi_M$ satisfies the hypothesis.  Its coefficient
vectors $(c^{(\nu)}_2,\ldots,c^{(\nu)}_M)$ converge to $(c_2,\ldots,c_M)$,
whose first nonzero entry has index $j_0-1$ in the new family, so by induction
$(f_\nu/\varphi_1)'=\sum_{j\ge2}c^{(\nu)}_j\chi_j$ has at most $j_0-2$ zeros
in $[n_\nu,\infty)$ for large $\nu$.  If $f_\nu$ had $j_0$ zeros there, so
would $f_\nu/\varphi_1$, and Rolle's theorem would give $j_0-1$ zeros of its
derivative between them.
\end{proof}

\begin{proposition}[Far zeros at repeated roots]\label{prop:farrep}
Let $r_1\le\cdots\le r_L$ with $r_1>1$, repetitions allowed, let
$z_1>\cdots>z_p$ be the distinct values of $1/r_i$ with multiplicities $m_i$,
and let $\psi_1,\ldots,\psi_L$ be the Hermite basis $d^az_i^d$,
$0\le a<m_i$, ordered by growth: $z$ descending, and $a$ descending within a
node.  For $0\le k\le L$ let
$\mathcal H_k=\operatorname{span}\{\psi_{k+1},\ldots,\psi_L\}$, and call
$v\in\mathcal H_k$ a certificate for $X$ on $\mathcal H_k$ if $v_0=1$ and
$v_x=0$ for $x\in X$.  Fix $0\le k<L$, and let $E_\nu$ be $k$-sets of positive
integers with $\min E_\nu\to\infty$.
\begin{enumerate}
\item[(a)] If $f_\nu\in\mathcal H_0$ vanishes on $E_\nu$ and its coefficient
vector has sup-norm $1$, its coefficients on $\psi_1,\ldots,\psi_k$ tend to
$0$.
\item[(b)] With $J^*=\{1,\ldots,k\}$ and
$\alpha_J(E)=\det(\psi_j(e))_{e\in E,\,j\in J}$, for large $\nu$ one has
$\alpha_{J^*}(E_\nu)\ne0$, and
$\eta_\nu=\max_{J\ne J^*}|\alpha_J(E_\nu)/\alpha_{J^*}(E_\nu)|\to0$.
\item[(c)] If $v^\infty$ is a certificate for a finite set $X_0$ on
$\mathcal H_k$, then for large $\nu$ there are certificates $v^{(\nu)}$ for
$X_0\cup E_\nu$ on $\mathcal H_0$ with $\Tail(v^{(\nu)}-v^\infty)\to0$.
\end{enumerate}
\end{proposition}

Removing the first $k$ functions removes whole blocks at the largest nodes
and, at the node where the cut falls, the top exponents; so $\mathcal H_k$ is
the Hermite space of the multiset $r_{k+1},\ldots,r_L$, and by the zero bound
for exponential polynomials~\cite[Part~Five, Ch.~1]{polya-szego} a nonzero
element of $\mathcal H_k$ has at most $L-k-1$ real zeros.  For distinct roots
$\mathcal H_k$ is the span of
$y_{k+1}^d,\ldots,y_L^d$, the case Sections~\ref{sec:value}
and~\ref{sec:subtwo} use.

\begin{proof}
(a) Otherwise, along a subsequence, the coefficient vectors converge to some
$c$ with a nonzero coordinate among the first $k$, so the least index $j_0$
with $c_{j_0}\ne0$ is at most $k$.  Each $\psi_j(x)=x^az_i^x$ lies in
$\mathcal C$, and consecutive ratios are $1/x$ within a node and
$x^{m_{i+1}-1}(z_{i+1}/z_i)^x$ across nodes, both tending to $0$.  By
Lemma~\ref{lem:infinity} with $n_\nu=\min E_\nu$, $f_\nu$ has at most
$j_0-1\le k-1$ zeros in $[\min E_\nu,\infty)$ for large $\nu$, but it vanishes
on the $k$ points of $E_\nu$.

(b) If $\alpha_{J^*}(E_\nu)=0$ along a subsequence, a combination of
$\psi_1,\ldots,\psi_k$ of sup-norm $1$ vanishes on $E_\nu$, against (a).
Otherwise let $T_\nu=\alpha_{J^*}(E_\nu)^{-1}(\psi_j(e))_{e\in E_\nu,\,j\le L}$,
a $k\times L$ matrix whose first $k$ columns form the identity; by Cramer's
rule $\alpha_J/\alpha_{J^*}$ is the minor of $T_\nu$ on the columns $J$.  For
$l>k$ the function $\psi_l-\sum_{j\le k}(T_\nu)_{jl}\psi_j$ vanishes on
$E_\nu$; its coefficient vector has sup-norm $\max(1,t)$ with
$t=\max_j|(T_\nu)_{jl}|$, so by (a) $t/\max(1,t)\to0$, that is $t\to0$.
Every $J\ne J^*$ contains some $l>k$, so its minor is a determinant with a
column tending to $0$ and the others bounded, and $\eta_\nu\to0$.

(c) Put $n=L-k$ and $s=n-1-|X_0|\ge0$, the zero bound in $\mathcal H_k$
giving $|X_0|\le n-1$, and fix $s$ positive integers $Q$ outside $X_0$.  For
large $\nu$ the $L$ points of $Y_\nu=\{0\}\cup X_0\cup Q\cup E_\nu$ are
distinct, so there is a unique $v^{(\nu)}\in\mathcal H_0$ with
$v^{(\nu)}_0=1$, $v^{(\nu)}|_{X_0}=0$, $v^{(\nu)}_q=v^\infty_q$ for $q\in Q$
and $v^{(\nu)}|_{E_\nu}=0$; it is a certificate for $X_0\cup E_\nu$, and
$v^\infty$ is the unique element of $\mathcal H_k$ meeting the same
conditions at the $n$ points $\{0\}\cup X_0\cup Q$.  Let $c^{(\nu)}$ be the
coefficient vector of $v^{(\nu)}$.  If $\|c^{(\nu)}\|_\infty\to\infty$ along
a subsequence, then $v^{(\nu)}/\|c^{(\nu)}\|_\infty$ has, by (a) and
compactness, a limit $g\in\mathcal H_k$ of sup-norm $1$ vanishing at the $n$
points $\{0\}\cup X_0\cup Q$, against the zero bound.  So the $c^{(\nu)}$ are
bounded, and also $\|c^{(\nu)}\|_\infty\ge1$ because $v^{(\nu)}_0=1$; by (a)
their first $k$ coordinates tend to $0$, so every limit point is the
coefficient vector of an element of $\mathcal H_k$ meeting the $n$ conditions
of $v^\infty$, that is, of $v^\infty$.  Hence $c^{(\nu)}\to c^\infty$, and
\[
  \Tail(v^{(\nu)}-v^\infty)\le\sum_{j=1}^L|c^{(\nu)}_j-c^\infty_j|
  \sum_{d\ge1}|\psi_j(d)|\longrightarrow0 .
\]
\end{proof}

For distinct roots, part (b) also has a direct proof with a rate.  Every
term of $\alpha_J(E)$, $J=\{j_1<\cdots<j_k\}$, is a product
$\prod_ty_{j_t}^{e_{\pi(t)}}$ for $E=\{e_1<\cdots<e_k\}$, and $\alpha_{J^*}$ is a
Vandermonde product times a Schur polynomial with nonnegative coefficients;
comparing the largest terms gives
$\eta\le(k!/c)(y_{k+1}/y_k)^{\min E}$ with $c=\prod_{a<b\le k}(1-y_b/y_a)$.
At repeated roots $c$ vanishes, and $\alpha_{J^*}$ can vanish outright: at
the roots $(3,3,4,4)$ with $k=3$ the leading functions
$d(\tfrac13)^d,(\tfrac13)^d,d(\tfrac14)^d$ have the combination
$\tfrac{2187}{16384}d(\tfrac13)^d-\tfrac{15309}{8192}(\tfrac13)^d+d(\tfrac14)^d$
vanishing at $d=6,7,8$.  This is why Proposition~\ref{prop:farrep} is proved
through Lemma~\ref{lem:infinity} rather than by estimating minors.  When $E$ has bounded diameter an
expansion of the confluent alternants gives
$\eta=O(1/\min E)$ if the cut after $\psi_k$ falls inside a multiplicity
block, and a geometric rate if it falls between blocks; for arbitrary $E$ no
rate is proved, and none is needed.

\begin{theorem}[Necessity of the partial-sum criterion]\label{thm:converse}
Let $2\le r_1\le\cdots\le r_L$, repetitions allowed, and $0\le u\le L-1$.  If
\eqref{eq:partial} fails, then
\[
  \inf_Qb_{u+1}\Bigl(Q\prod_{i=1}^L(x-r_i)\Bigr)>\prod_{i=u+1}^L(r_i-1)
\]
over monic $Q\in\R[x]$.  With Corollary~\ref{cor:charrep},
\eqref{eq:order} is the infimum exactly when \eqref{eq:partial} holds.
\end{theorem}

\begin{proof}
Use the notation of Proposition~\ref{prop:farrep}.  If $F$ is a monic
multiple of $P=\prod_i(x-r_i)$ of degree $D$ and $v\in\mathcal H_0$, then
$\sum_{j\le D}f_jv_{D-j}=0$: for $\psi=d^az^d$ with $z=1/r$, expanding
$(D-j)^a$ in powers of $j$ writes $\sum_jf_j(D-j)^ar^{j-D}$ as a combination
of the values $\bigl((x\,d/dx)^bF\bigr)(r)$, $b\le a$, which vanish because
$r$ is a root of multiplicity greater than $a$.  So \eqref{eq:weak} holds
for every certificate on $\mathcal H_0$, and, as in the distinct case, it
suffices to find $\delta>0$ such that every $X\subset\N_{>0}$ of size $u$ has
a certificate on $\mathcal H_0$ with $\Tail\le B-\delta$.  Suppose not, and
take $X_\nu$ whose certificates on $\mathcal H_0$ all have
$\Tail\ge B-1/\nu$.  Passing to a subsequence, $X_\nu=X_0\cup E_\nu$ with
$X_0$ fixed, $|E_\nu|=k$ and $\min E_\nu\to\infty$.  By
Proposition~\ref{prop:farrep}(c), a certificate $v^\infty$ for $X_0$ on
$\mathcal H_k$ with $\Tail(v^\infty)<B$ yields certificates for $X_\nu$
with tails tending to $\Tail(v^\infty)$, a contradiction; so it suffices to
produce one.  Put
$n=L-k$, $m=L-u$ and $u_0=|X_0|=u-k$; recall that $\mathcal H_k$ is the
Hermite space of $r_{k+1},\ldots,r_L$.

If $u_0=0$, $\mathcal H_u$ is the Hermite space of the roots of $P_u$, and
since \eqref{eq:partial} fails, Lemma~\ref{lem:improve} (with repetitions)
for $Q=P_u$ and $X_0=\varnothing$ gives $\Tail(v^\infty)<1/|P_u(1)|=B$.

If $u_0\ge1$, take the certificate of Lemma~3.1's proof on $\mathcal H_k$:
zero set $Z$ consisting of $X_0$ and the first $m-1$ positive integers
outside it, so $|Z|=n-1$ and its leading run is $f\ge m-1$.
Proposition~\ref{prop:confluent}, with nodes at most $1/2$, bounds $\Tail$ by
the product of $1/(r_i-1)$ over the $f+1$ largest of $r_{k+1},\ldots,r_L$,
strictly unless $Z=\{1,\ldots,n-1\}$ by Lemma~\ref{lem:confstrict}.  Every
factor is at most $1$ and $f+1\ge m$, so $\Tail\le B$, strictly unless
$Z=\{1,\ldots,n-1\}$ and $r_{k+1}=\cdots=r_u=2$.  In that case
$X_0\subseteq\{1,\ldots,n-1\}$ and $Q=\prod_{i>k}(x-r_i)=(x-2)^{u_0}P_u$
has $|Q(1)|=|P_u(1)|=1/B$, so Lemma~\ref{lem:improve} (with repetitions) for
$Q$ and $X_0$ finishes the proof once some $j\in\{1,\ldots,n-1\}\setminus X_0$
has $|S_j(Q)|>|Q(1)|$.

For a monic $R$ with all roots above $1$ let
$\mathcal A(R)=\{1\le j\le\deg R-1:|S_j(R)|>|R(1)|\}$.  The coefficients of
$(x-2)R$ read from the top are $c_l-2c_{l-1}$, so
$S_j((x-2)R)=S_j(R)-2S_{j-1}(R)$, and by Lemma~\ref{lem:signs}
$(-1)^jS_j(R)>0$ for $0\le j\le\deg R$; hence
$|S_j((x-2)R)|=|S_j(R)|+2|S_{j-1}(R)|$ for $1\le j\le\deg R$, while
$|((x-2)R)(1)|=|R(1)|$.  Thus $j\in\mathcal A(R)$ puts both $j$ and $j+1$ in
$\mathcal A((x-2)R)$.  Since \eqref{eq:partial} fails, $\mathcal A(P_u)$
contains some $i\le m-1$, and $u_0$ applications give
$\{i,i+1,\ldots,i+u_0\}\subseteq\mathcal A(Q)\subseteq\{1,\ldots,n-1\}$.
These are $u_0+1$ integers and $X_0$ has only $u_0$ elements.
\end{proof}

For distinct roots the equality case $Z=\{1,\ldots,n-1\}$,
$r_{k+1}=\cdots=r_u=2$ arises only for $u=1$, $r_1=2$ and $k=0$; with
repetitions it also covers $r_1=\cdots=r_u=2$ for any $u$ and any $k<u$.

For distinct roots the proof also shows how a multiple can approach
\eqref{eq:order}.  Outside
the case $u=1$, $r_1=2$, every placement with an exempt coefficient at
bounded distance from the top has a certificate with tail below $B$, so
along any sequence of monic multiples with $b_{u+1}$ tending to the bound all
$u$ exempt coefficients must move away from the top, as the construction of
Proposition~\ref{prop:partial} moves them to the bottom.  When
\eqref{eq:partial} fails, Section~\ref{sec:value} computes the infimum.

\section{The infimum when the criterion fails}\label{sec:value}

Keep the notation of the necessity argument in
Section~\ref{sec:criterion}: $r_1<\cdots<r_L$ are distinct, $y_i=1/r_i$, and $m=L-u$.  For finite $X\subset\N_{>0}$ and
$0\le k<L$ put
\[
  \tau_k(X)=\inf\bigl\{\Tail(v):v\text{ a certificate for }X\text{ on }
  y_{k+1},\ldots,y_L\bigr\},
  \quad \tau(X)=\tau_0(X),\quad \tau^*=\tau_u(\varnothing).
\]
A certificate on $c$ nodes is \emph{full} if it has $c-1$ prescribed zeros.
It is then unique, and by the zero bound it has no other zeros and keeps one
sign past the largest.  By \eqref{eq:weak} every monic multiple $F$ has
$b_{u+1}(F)\ge1/\tau(X_F)$, where $X_F$ holds the distances from the top of
its $u$ largest nonleading coefficients.  So the infimum is set by the worst
placement.  The first result says that it is exactly that, and that exempt
coefficients escaping to the bottom cost $\tau^*$, the best tail on the
surviving roots alone.

\begin{theorem}[The infimum as a placement game]\label{thm:value}
Let $1<r_1<\cdots<r_L$ and $0\le u\le L-1$.  Then
\[
  \inf_Qb_{u+1}\Bigl(Q\prod_{i=1}^L(x-r_i)\Bigr)=\frac1T,
  \qquad T=\sup_{|X|=u}\tau(X)\ge\tau^*,
\]
over monic $Q\in\R[x]$.
\end{theorem}

\begin{proof}
The lower bound $b_{u+1}(F)\ge1/\tau(X_F)\ge1/T$ is \eqref{eq:weak}.

For the upper bound fix $X$ with $|X|=u$ and $D\ge\max(L,\max X)$, and let
$\tau_D(X)$ be the least $\sum_{d=1}^D|v_d|$ over certificates $v$ for $X$.
It is positive, since a certificate vanishing at $1,\ldots,D$ would vanish
identically.  Minimizing the largest $|f_j|$ over the positions $j<D$ with
$D-j\notin X$, subject to $F=x^D+\sum_{j<D}f_jx^j$ vanishing at every $r_i$,
is a feasible linear program.  Write $\mu_i$ for the multiplier of
$F(r_i)=0$, and put $\lambda_i=-\mu_ir_i^D$ and $v_d=\sum_i\lambda_iy_i^d$.
The coefficient of $f_j$ in the Lagrangian is then $-v_{D-j}$, and the dual
program is to maximize $v_0$ subject to $v_d=0$ for $d\in X$ and
$\sum_{d\le D,\,d\notin X}|v_d|\le1$.  Its value is $1/\tau_D(X)$.  By
linear-programming duality some monic multiple of degree $D$ has every
coefficient off the $u$ positions of $X$ at most $1/\tau_D(X)$ in magnitude,
so its $b_{u+1}$ is at most $1/\tau_D(X)$.

Two limits finish the proof.  First, $\tau_D(X)$ is nondecreasing in $D$ and
at most $\tau(X)$.  Take minimizers $v^{(D)}=\sum_i\lambda^{(D)}_iy_i^d$.
Their values at $L$ fixed points of $\{0,1,\ldots\}\setminus X$, one of them
$0$, are bounded by $\max(1,\tau(X))$.  The matrix $(y_i^d)$ at $L$ distinct
points is nonsingular by the zero bound, so the $\lambda^{(D)}$ are bounded.
A subsequence converges to a certificate $v$ for $X$ with
$\sum_{d\le N}|v_d|\le\lim_D\tau_D(X)$ for every $N$, whence
$\tau(X)\le\Tail(v)\le\lim_D\tau_D(X)$.  So the infimum is at most $1/\tau(X)$
for every $X$, that is, at most $1/T$.

Second, put $E_\nu=\{\nu+1,\ldots,\nu+u\}$ and suppose
$\tau(E_\nu)\le\tau^*-\varepsilon$ along a subsequence.  Take certificates
$v^\nu=\sum_i\lambda^\nu_iy_i^d$ for $E_\nu$ with
$\Tail(v^\nu)\le\tau^*-\varepsilon/2$.  For $\nu\ge L$ their values at
$0,\ldots,L-1$ are bounded, so we may assume $\lambda^\nu\to\lambda$.  The $u$
equations $v^\nu_e=0$, $e\in E_\nu$, read
$A_\nu\lambda^\nu_{\le u}=-R_\nu\lambda^\nu_{>u}$ with
$A_\nu=(y_j^e)_{e\in E_\nu,\,j\le u}$ and
$R_\nu=(y_j^e)_{e\in E_\nu,\,j>u}$.  By Cramer's rule each entry of
$A_\nu^{-1}R_\nu$ is, up to sign, a ratio $\alpha_J/\alpha_{J^*}$ of
Proposition~\ref{prop:farrep}(b), whose basis is $y_1^d,\ldots,y_L^d$ for
distinct roots, with $k=u$ and $J\ne J^*$.  So it is at most $\eta_\nu\to0$, and $\lambda_{\le u}=0$.  Then
$v=\sum_{i>u}\lambda_iy_i^d$ is a certificate on the surviving nodes, and as
before $\Tail(v)\le\tau^*-\varepsilon/2$, against the definition of $\tau^*$.
Hence $\liminf_\nu\tau(E_\nu)\ge\tau^*$, and $T\ge\tau^*$.
\end{proof}

Proposition~\ref{prop:farrep}(c) with $k=u$ and $X_0=\varnothing$ supplies the
reverse inequality, so $\tau(E_\nu)\to\tau^*$:
the escaping placement costs exactly $\tau^*$.  There is also a dual side.  If
$v=\sum_{i>u}\lambda_iy_i^d$ is a certificate on the surviving nodes and $w$
is a bounded sequence with $w_0=1$ whose series $W(y)=\sum_dw_dy^d$ vanishes
at $y_{u+1},\ldots,y_L$, then $0=\sum_i\lambda_iW(y_i)=\sum_dw_dv_d$, so
\begin{equation}\label{eq:dualpair}
  1\le\sup_{d\ge1}|w_d|\cdot\Tail(v).
\end{equation}
Such $W$ are the $\widetilde P_uH$ with $H(0)=1$.  Taking
$H=\widetilde M/(1-y)$ for a monic $M\in\R[x]$ with reverse $\widetilde M$
makes $w_d$ the partial sums $S_d(P_uM)$, as in \eqref{eq:wgen}.  So
\begin{equation}\label{eq:multiplier}
  \inf_Qb_{u+1}\Bigl(Q\prod_{i=1}^L(x-r_i)\Bigr)\le\frac1{\tau^*}
  \le\max_{j\ge1}\bigl|S_j(P_uM)\bigr|
  \qquad\text{for every monic }M\in\R[x].
\end{equation}
At $M=1$ this is the construction of Proposition~\ref{prop:partial}.  A
witness $W$ with $\sup_{d\ge1}|w_d|=1/\Tail(v)$ proves that $v$ attains
$\tau^*$.  Such a witness always exists.

\begin{proposition}[Attainment with a multiplier witness]\label{prop:attain}
Write $m=L-u$.  Then $\tau^*$ is attained by a full certificate $v$ on
$y_{u+1},\ldots,y_L$.  At any such $v$ there are reals $s_d$ with
$|s_d|\le1$, $s_d=\sgn(v_d)$ whenever $v_d\ne0$, and
$\sum_{d\ge1}s_dy_i^d=\tau^*$ for $u<i\le L$; then $w_0=1$,
$w_d=-s_d/\tau^*$ is a witness as in \eqref{eq:dualpair} attaining equality,
and it is $\widetilde P_u\widetilde M/(1-y)$ for a monic $M$, so
\eqref{eq:multiplier} is an equality at that $M$.
\end{proposition}

\begin{proof}
A certificate for $\varnothing$ on the surviving nodes is a
$\lambda\in\R^m$ with $\sum_i\lambda_i=1$, writing
$v_d(\lambda)=\sum_{i>u}\lambda_iy_i^d$; put $f(\lambda)=\Tail(v(\lambda))$.

\emph{The infimum is attained.}  Since $|v_d|\le\sum_i|\lambda_i|y_i^d$ and
the nodes lie in $(0,1)$, $f$ is finite, and it is homogeneous and
subadditive.  If $f(\lambda)=0$ then $v_d=0$ for $1\le d\le m$, a Vandermonde
system in the $\lambda_i$ with distinct nonzero nodes, so $\lambda=0$.  Thus
$f$ is a norm, hence continuous and coercive, and it attains its minimum on
the closed nonempty affine set $\{\sum_i\lambda_i=1\}$; that minimum is
$\tau^*$.

\emph{Some minimiser is full.}  Let $\lambda$ minimise, let
$Z=\{d\ge1:v_d=0\}$, and suppose $|Z|\le m-2$.  The conditions
$\sum_ih_i=0$ and $v_d(h)=0$ for $d\in Z$ are at most $m-1$ linear
conditions on $h\in\R^m$, so some $h\ne0$ meets them.  As $v_d(h)=0$ on $Z$,
\[
  \varphi(t)=f(\lambda+th)=\sum_{d\notin Z}\bigl|v_d+tv_d(h)\bigr|,
\]
and each term is differentiable at $t=0$ with difference quotients dominated
by the summable $|v_d(h)|$, so $\varphi$ is differentiable there and
minimality gives $\varphi'(0)=\sum_{d\notin Z}\sgn(v_d)v_d(h)=0$.
Subtracting that vanishing linear term,
\[
  \varphi(t)-\varphi(0)=\sum_{d\notin Z}
  \bigl(|v_d+tv_d(h)|-|v_d|-t\sgn(v_d)v_d(h)\bigr),
\]
all of whose terms are nonnegative, the $d$-th vanishing exactly while
$v_d+tv_d(h)$ keeps the sign of $v_d$, that is for $t$ between $0$ and
$t_d=-v_d/v_d(h)$ (no condition when $v_d(h)=0$).  Neither $v(\lambda)$ nor
$v(h)$ vanishes identically, so by the zero bound each has constant nonzero
sign past its largest zero, and hence $t_d$ has one sign for all large $d$.
On the other side lie only finitely many $t_d$, all nonzero, so the nearest
of them is at a positive distance $t^*$, and $t^*$ is finite because
otherwise $\varphi$ would be constant on a half-line, against coercivity.
On the interval from $0$ to $t^*$ every term above vanishes, so
$\varphi(t^*)=\varphi(0)$: the certificate $\lambda+t^*h$ still minimises,
keeps every zero of $Z$, and has gained one.  The zero bound caps the zeros
at $m-1$, so iterating reaches a full minimiser.

\emph{The multipliers.}  Let $\lambda$ be a full minimiser with zero set
$Z$, $|Z|=m-1$, and set $s_d=\sgn(v_d)$ for $d\notin Z$.  A nontrivial real
combination of the $m$ distinct powers $y^d$, $d\in\{0\}\cup Z$, has at most
$m-1$ positive zeros, so the matrix $(y_i^d)_{i>u,\,d\in\{0\}\cup Z}$ is
nonsingular.  Hence there are unique reals $(s_d)_{d\in Z}$ and $\theta$
with $\sum_{d\ge1}s_dy_i^d=\theta$ for every surviving $i$, the series
converging because the $s_d$ are bounded.  Pairing with any $h$ with
$\sum_ih_i=0$ gives $\sum_{d\ge1}s_dv_d(h)=\theta\sum_ih_i=0$, while
minimality gives
\[
  0\le f'(\lambda;h)=\sum_{d\notin Z}\sgn(v_d)v_d(h)+\sum_{d\in Z}|v_d(h)|,
\]
so subtracting the pairing leaves
$0\le\sum_{d\in Z}\bigl(|v_d(h)|-s_dv_d(h)\bigr)$.  By the same
nonsingularity $h\mapsto(v_d(h))_{d\in Z}$ carries $\{\sum_ih_i=0\}$ onto
$\R^{m-1}$, so taking the preimages of $\pm e_d$ gives $|s_d|\le1$ on $Z$.
Finally $\sum_{d\ge1}s_dv_d=\sum_i\lambda_i\sum_{d\ge1}s_dy_i^d=\theta$,
while $v_d=0$ on $Z$ makes the left side $\sum_{d\notin Z}|v_d|=\Tail(v)$,
so $\theta=\tau^*$.

\emph{The witness is a monic multiplier.}  Put $w_0=1$ and
$w_d=-s_d/\tau^*$.  Then $W(y_i)=1-\tau^*/\tau^*=0$ at every surviving node,
and $\sup_{d\ge1}|w_d|=1/\tau^*$ since $|s_d|=1$ off $Z$ and $|s_d|\le1$ on
it, so \eqref{eq:dualpair} is an equality.  Past $g=\max Z$ the sign of
$v_d$ is constant by the zero bound, so $w_d$ is constant there and
$(1-y)W(y)$ is a polynomial of degree at most $g+1$, vanishing at the $m$
surviving nodes and equal to $1$ at $0$.  Dividing out those nodes writes it
as $\widetilde P_uN$ with $N(0)=1$, and the reverse of $N$ is a monic $M$ of
degree at most $g+1-m$ with $N=\widetilde M$.
\end{proof}

The dual side also settles sufficiency of the criterion at repeated roots,
where the construction of Proposition~\ref{prop:partial} is not available.

\begin{corollary}[The criterion at repeated roots]\label{cor:charrep}
Let $2\le r_1\le\cdots\le r_L$, repetitions allowed, and $0\le u\le L-1$.
Then \eqref{eq:order} is the infimum of $b_{u+1}$ over monic multiples of
$\prod_i(x-r_i)$ exactly when \eqref{eq:partial} holds; no alignment
$r_u<r_{u+1}$ is needed.
\end{corollary}

\begin{proof}
Necessity is Theorem~\ref{thm:converse}.  For sufficiency use
Theorem~\ref{thm:value}, which carries to repeated roots: in its first half
the constraints $F(r_i)=0$ become $((x\,d/dx)^aF)(r_i)=0$, $a<m_i$, the dual
variables again assemble into an element of $\mathcal H_0$ because
$d\mapsto(D-d)^ay_i^{d}$ and $d^ay_i^d$ span the same space, and the
compactness step uses unisolvence of $\mathcal H_0$ at $L$ distinct points;
in its second half the entries of $A_\nu^{-1}R_\nu$ are the ratios of
Proposition~\ref{prop:farrep}(b), which tend to $0$ for every sequence with
$\min E_\nu\to\infty$.  So the infimum is $1/T$ with $T\ge\tau^*$, where
$\tau^*$ is the least tail of a certificate on $\mathcal H_u$.  The dual
pairing \eqref{eq:dualpair} holds on $\mathcal H_u$, since
$\sum_dw_dd^ay^d=((y\,d/dy)^aW)(y)$ vanishes at $y=1/r_i$ for $a<m_i$ when
$\widetilde P_u$ divides $W$; with $M=1$ in \eqref{eq:multiplier} it gives
$1/\tau^*\le\max_{j\ge1}|S_j(P_u)|$, which is $|P_u(1)|=1/B$ under
\eqref{eq:partial}.  Hence $T\ge B$, while Theorem~3.2 gives $T\le B$.
\end{proof}

\paragraph{Hermite spaces.}  For real roots
$1<r_1\le\cdots\le r_L$, repetitions allowed, $\mathcal H_k$ is the Hermite
space of the multiset $r_{k+1},\ldots,r_L$ as in
Proposition~\ref{prop:farrep}: the span of $d\mapsto d^az^d$, $z=1/r$ a
distinct value among them and $a$ below its multiplicity there, ordered by
growth.  Its dimension is $L-k$, and by the P\'olya--Szeg\H o zero bound a
nonzero element has at most $L-k-1$ real zeros counted with multiplicity.
A \emph{certificate for $X$ on $\mathcal H_k$} is a $v\in\mathcal H_k$ with
$v_0=1$ and $v|_X=0$; it is \emph{full} if it is prescribed $L-k-1$ zeros.  A
full certificate is unique (unisolvence at distinct points), has no other
real zero, changes sign at each prescribed zero and keeps one sign past the
largest.  For finite $X\subset\N_{>0}$ put
\[
  \tau_k(X)=\inf\{\Tail(v):v\text{ a certificate for }X\text{ on }\mathcal H_k\},
  \qquad \tau(X)=\tau_0(X),\qquad \tau^*=\tau_u(\varnothing).
\]
For distinct roots these are the quantities of Section~\ref{sec:value}.  The top basis function $\psi_{k+1}$ of $\mathcal H_k$ is
$d^{e-1}y^d$, where $y=1/r_{k+1}$ is the largest node of $\mathcal H_k$ and
$e$ its multiplicity there; $\mathcal H_{k+1}$ is the span of the others.

\begin{theorem}[The placement game at repeated roots]\label{thm:valuerep}
Let $1<r_1\le\cdots\le r_L$, repetitions allowed, and $0\le u\le L-1$.  Then
\[
  \inf_Qb_{u+1}\Bigl(Q\prod_{i=1}^L(x-r_i)\Bigr)=\frac1T,
  \qquad T=\sup_{|X|=u}\tau(X)\ge\tau^*,
\]
over monic $Q\in\R[x]$.
\end{theorem}

\begin{proof}
The lower bound is \eqref{eq:weak}, which holds for every certificate on
$\mathcal H_0$ by the first paragraph of the proof of
Theorem~\ref{thm:converse}.

For the upper bound fix $X$ with $|X|=u$ and $D\ge\max(L,\max X)$, and let
$\tau_D(X)$ be the least $\sum_{d=1}^D|v_d|$ over certificates $v$ for $X$
on $\mathcal H_0$; it is positive because a nonzero element of
$\mathcal H_0$ has at most $L-1$ zeros.  Write $m_i$ for the multiplicity of
$r_i$ and $\theta=x\,d/dx$.  Since $r_i\ne0$, $F=x^D+\sum_{j<D}f_jx^j$ is a
multiple of $P$ exactly when
$(\theta^aF)(r_i)=\sum_{j\le D}f_jj^ar_i^j=0$ for $a<m_i$.  Minimising the
largest $|f_j|$ over the positions $j<D$ with $D-j\notin X$ subject to these
constraints is a feasible linear program ($x^{D-L}P$ is feasible).  With
multipliers $\mu_{i,a}$ put
$v_d=-\sum_{i,a}\mu_{i,a}r_i^D(D-d)^ay_i^d$; then the coefficient of $f_j$ in
the Lagrangian is $-v_{D-j}$, and as $\{(D-d)^a\}_{a<m_i}$ and
$\{d^a\}_{a<m_i}$ span the same polynomials, $\mu\mapsto v$ is a linear
bijection onto $\mathcal H_0$.  The dual program is to maximise $v_0$
subject to $v|_X=0$ and $\sum_{d\le D,\,d\notin X}|v_d|\le1$, with value
$1/\tau_D(X)$, so some monic multiple of degree $D$ has
$b_{u+1}\le1/\tau_D(X)$.  The limit $\tau_D(X)\uparrow\tau(X)$ is the
compactness step of the proof of Theorem~\ref{thm:value}, using unisolvence
of $\mathcal H_0$ at $L$ distinct points; so the infimum is at most $1/T$.

For $T\ge\tau^*$ put $E_\nu=\{\nu+1,\ldots,\nu+u\}$ and suppose
$\tau(E_\nu)\le\tau^*-\varepsilon$ along a subsequence; take certificates
$v^\nu$ for $E_\nu$ on $\mathcal H_0$ with
$\Tail(v^\nu)\le\tau^*-\varepsilon/2$.  For $\nu\ge L$ their values at
$0,\ldots,L-1$ are bounded, so their coefficient vectors $c^\nu$ are bounded,
and $\|c^\nu\|_\infty$ is bounded below because $v^\nu_0=1$.
Proposition~\ref{prop:farrep}(a) with $k=u$, applied to
$v^\nu/\|c^\nu\|_\infty$, sends the coordinates of $c^\nu$ on
$\psi_1,\ldots,\psi_u$ to $0$.  A convergent subsequence therefore has a
limit $v\in\mathcal H_u$ with $v_0=1$, and by Fatou's lemma
$\Tail(v)\le\tau^*-\varepsilon/2$, against the definition of $\tau^*$.
\end{proof}

\begin{proposition}[Attainment at repeated roots]\label{prop:attainrep}
In the setting of Theorem~\ref{thm:valuerep}, $\tau^*$ is attained by a full
certificate on $\mathcal H_u$, and some monic $M\in\R[x]$ has
$\max_{j\ge1}|S_j(P_uM)|=1/\tau^*$; moreover
$\max_{j\ge1}|S_j(P_uM)|\ge1/\tau^*$ for every monic $M$.
\end{proposition}

\begin{proof}
Run the proof of Proposition~\ref{prop:attain} on $\mathcal H_u$, with the
coefficient vector $\lambda$ in the Hermite basis, the constraint $v_0=1$ in
place of $\sum_i\lambda_i=1$, and directions $h$ with $h_0=0$.  $\Tail$ is a
norm on $\mathcal H_u$ (an element vanishing at $1,\ldots,m$ vanishes), a
nonzero element has finitely many real zeros and so one sign for large $d$,
and the zero bound caps the zeros at $m-1$; so the first two steps hold
verbatim.  At a full minimiser with zero set $Z$, unisolvence at
$\{0\}\cup Z$ gives unique $(s_d)_{d\in Z}$ and $\theta$ with
$\sum_{d\ge1}s_dh_d=\theta h_0$ for all $h\in\mathcal H_u$ (with
$s_d=\sgn v_d$ off $Z$), and makes $h\mapsto(h_d)_{d\in Z}$ map
$\{h_0=0\}$ onto $\R^{m-1}$; so $|s_d|\le1$ and $\theta=\tau^*$ as there.
Put $w_0=1$, $w_d=-s_d/\tau^*$.  Then $\sum_dw_dh_d=0$ on $\mathcal H_u$;
with $h=d^ay_i^d$ this reads $((y\,d/dy)^aW)(y_i)=0$ for $a$ below the
multiplicity of $y_i$, so $W$ vanishes there to that order.  Past $\max Z$
the sign of $v$ is constant, $(1-y)W$ is a polynomial divisible by
$\widetilde P_u=\prod_{i>u}(1-r_iy)$ with value $1$ at $0$, and the quotient
is $\widetilde M$ for a monic $M$, whose partial sums $S_j(P_uM)$ are the
$w_j$.  The last claim is \eqref{eq:dualpair} on $\mathcal H_u$ (proof of
Corollary~\ref{cor:charrep}).
\end{proof}

\subsection*{The escaping placement is extremal}

Whether $T=\tau^*$, so that the escaping placement is extremal, comes down to
covering the finite placements.  The deletion step of Theorem~2.2 lifts a
certificate on fewer nodes to one on all of them.

\begin{theorem}[Lifting]\label{thm:lift}
Let $2\le r_1<\cdots<r_L$ and $0\le k\le u$, and let $v$ be a full
certificate on $y_{k+1},\ldots,y_L$ with zero set $Z$ and $g=\max Z$ ($g=0$
if $Z$ is empty).  Then $\tau(X)\le\Tail(v)$ for every $X$ of size $u$ with
$X\subseteq Z\cup(g,\infty)$ and $|X\setminus Z|\le k$.  The inequality is
strict if $k\ge1$ and the padded set $Z^+$ below is not $\{1,\ldots,L-1\}$.
\end{theorem}

\begin{proof}
Add to $Z\cup X$ integers above $\max(Z\cup X)$ until it has $L-1$ elements,
and call the result $Z^+$.  It exceeds $Z$ by $k$ elements, all above $g$.
The full certificate $v^+$ on $y_1,\ldots,y_L$ with zero set $Z^+$ is a
certificate for $X$.  Delete the largest zero of $Z^+$ together with the
largest node $y_1$.  The nodes are at most $1/2$, so the deletion step in the
proof of Theorem~2.2 bounds $\Tail(v^+)$ by the tail of the full certificate
on $y_2,\ldots,y_L$ with the remaining zeros, strictly unless
$Z^+=\{1,\ldots,L-1\}$.  In that case Lemma~2.1 gives the two tails as
products that differ by the factor $1/(r_1-1)\le1$.  Repeat $k$ times.  Each
removed zero is the largest remaining, because the added ones all exceed $g$,
and the last step lands on $v$.
\end{proof}

The hypothesis $X\subseteq Z\cup(g,\infty)$ matters: inserting zeros one at a
time, each insertion paid for by one node, fails outright.  The full
certificate on $(\tfrac13,\tfrac15,\tfrac17)$ with zero set $\{1,3\}$ has
tail $\tfrac{31}{1704}$; adding the node $\tfrac12$ and the zero $2$ gives
the full certificate on all four nodes with the consecutive zero set
$\{1,2,3\}$, of tail $\tfrac1{48}$, so the tail rose although a node was
gained---Theorem~2.2's equality case absorbed it.  Barring consecutive zero
sets does not repair the step.  At $(2,\tfrac94,\tfrac52,3)$ the full
certificate on $(\tfrac49,\tfrac25,\tfrac13)$ with zero set $\{1,4\}$ has
tail $\tfrac{11608}{84555}=0.13728\ldots$, and adding the node $\tfrac12$ and
the zero $2$ gives the zero set $\{1,2,4\}$ and tail
$\tfrac{108}{755}=0.14304\ldots$.

\begin{corollary}[One exempt coefficient]\label{cor:uone}
Let $2\le r_1<\cdots<r_L$ and $u=1$, and suppose $\tau^*$ is attained by a
full certificate on $y_2,\ldots,y_L$ with zero set $Z^*$ and $g=\max Z^*$.
Then
\[
  \inf_Qb_2\Bigl(Q\prod_{i=1}^L(x-r_i)\Bigr)
  =\frac1{\max\bigl(\tau^*,\ \max_{w\in G}\tau(\{w\})\bigr)},
  \qquad G=\{1,\ldots,g-1\}\setminus Z^*.
\]
\end{corollary}

\begin{proof}
Here $T=\sup_w\tau(\{w\})$.  If $w\notin G$, then $w\in Z^*$ or $w>g$, and
Theorem~\ref{thm:lift} with $k=1$ gives $\tau(\{w\})\le\tau^*$.  Since
$T\ge\tau^*$ by Theorem~\ref{thm:value}, the maximum is $T$.
\end{proof}

By Theorem~\ref{thm:extremal} the maximum is always $\tau^*$, so the
infimum of $b_2$ is $1/\tau^*$.

For general $u$ the question $T=\tau^*$ reduces to a single step, in which
one node is added and one zero is deleted, the largest.

\begin{proposition}[Step reduction]\label{prop:step}
Let $1<r_1<\cdots<r_L$ and $0\le u\le L-1$.  Suppose that for every
$0\le k<u$ and every set $X$ of $u-k$ positive integers
\begin{equation}\label{eq:step}
  \tau_k(X)\le\tau_{k+1}\bigl(X\setminus\{\max X\}\bigr).
\end{equation}
Then $T=\tau^*$, and the infimum of $b_{u+1}$ is $1/\tau^*$.
\end{proposition}

\begin{proof}
Chaining \eqref{eq:step} from $k=0$ upwards, each step deleting the largest
remaining element, gives
$\tau(X)=\tau_0(X)\le\tau_u(\varnothing)=\tau^*$ for every $X$ of size $u$, so
$T\le\tau^*$; Theorem~\ref{thm:value} gives the reverse.
\end{proof}

Deleting $\min X$ instead is false: at $(2,3,4,5)$ with $u=2$ and
$X=\{2,3\}$, $\tau_0(\{2,3\})=\tfrac1{24}$ against
$\tau_1(\{3\})=\tfrac{35}{1128}$, a ratio of $1.34$, while
$\tau_1(\{2\})=\tfrac1{24}$ meets it.  Theorem~\ref{thm:lift} with $k=1$
proves \eqref{eq:step} whenever the placement and the certificate it is lifted
from are aligned.

\begin{proposition}[The step in the aligned cases]\label{prop:stepeasy}
Let $2\le r_1<\cdots<r_L$, let $X$ be finite and nonempty with $g=\max X$ and
$X'=X\setminus\{g\}$, and let $v^-$ be a full certificate on
$y_{k+2},\ldots,y_L$ whose zero set $Z^-$ contains $X'$.  If $g\in Z^-$ or
$g>\max Z^-$, then some certificate for $X$ on $y_{k+1},\ldots,y_L$ has tail at
most $\Tail(v^-)$.
\end{proposition}

\begin{proof}
Put $Z^+=Z^-\cup\{g\}$ in the second case, and $Z^+=Z^-\cup\{N\}$ with any
$N>\max Z^-$ in the first.  Either way $Z^+\supseteq X$ has $L-k-1$ elements
and its largest exceeds all the others, so the full certificate $v^+$ on
$y_{k+1},\ldots,y_L$ with zero set $Z^+$ is a certificate for $X$; deleting
that largest zero together with the largest node $y_{k+1}$ is the deletion step
in the proof of Theorem~2.2, which bounds $\Tail(v^+)$ by $\Tail(v^-)$ because
the nodes are at most $1/2$.
\end{proof}

So \eqref{eq:step} is left open only when a minimizing zero set $Z^-$ for
$X\setminus\{\max X\}$ straddles $g=\max X$: $g\notin Z^-$ and
$g<\max Z^-$.  In that case the rule is not to displace one zero but to shift
every zero above $g$ up by one.  Write $Z^-=B_0\cup A$ with
$B_0=Z^-\cap(0,g)$ and $A=Z^-\cap(g,\infty)$, and $A+1=\{a+1:a\in A\}$.
Write $A$ as the disjoint union of its maximal runs of consecutive integers
$[s_1,e_1],\ldots,[s_p,e_p]$, so that
$A+1=\bigcup_j[s_j+1,e_j+1]$.  The neighbour of $B_0\cup A$ belonging to the
$j$th run is the zero set obtained by moving its first element past its last,
\[
  Z_j=(B_0\cup A)\setminus\{s_j\}\cup\{e_j+1\},
\]
which shifts that run up by one; for a run of length one it is the move
$a\to a+1$.  Shifting every run at once and adjoining $g$ gives the zero set
$B_0\cup\{g\}\cup(A+1)$ of the step.

\begin{proposition}[The run chord]\label{prop:runchord}
Let $y_{k+1}\le\tfrac12$, let $g\ge1$, let $B_0\subseteq\{1,\ldots,g-1\}$ and
let $A$ be a nonempty set of integers above $g$ with
$|B_0|+|A|=L-k-2$, with runs $[s_j,e_j]$ and neighbours $Z_j$ as above.  Let
$F$ and $F_j$ be the full certificates on $y_{k+2},\ldots,y_L$ with zero sets
$Z=B_0\cup A$ and $Z_j$, and $w$ the full certificate on $y_{k+1},\ldots,y_L$
with zero set $W=B_0\cup\{g\}\cup(A+1)$.  Then there are
$\lambda_1,\ldots,\lambda_p,\mu>0$ with $\sum_j\lambda_j+\mu=1$ such that
\begin{equation}\label{eq:runchord}
  \Tail(F)\ \ge\ \sum_{j=1}^p\lambda_j\Tail(F_j)+\mu\,\Tail(w),
\end{equation}
with equality if and only if $y_{k+1}=\tfrac12$ and $B_0=\{1,\ldots,g-1\}$.
\end{proposition}

\begin{proof}
\emph{Signs.}  For an integer $d\ge0$ put
$\varphi(d)=(-1)^{|Z\cap(0,d)|}$.  Each of $F$, $F_j$, $w$ is full, so by the
zero bound it has exactly its prescribed zeros, all simple, and it is positive
at $0$; its sign at a point $d$ outside its zero set is $-1$ to the number of
its zeros below $d$.  Hence $F(d)=\varphi(d)|F(d)|$.  For an integer $d$,
$|Z_j\cap(0,d)|-|Z\cap(0,d)|=[e_j+1<d]-[s_j<d]$, which vanishes unless
$s_j<d\le e_j+1$, and those $d$ lie in $Z_j$; so $F_j(d)=\varphi(d)|F_j(d)|$
for every integer $d\ge0$.  Since $|(A+1)\cap(0,d)|=|A\cap(0,d-1)|$,
\[
  |W\cap(0,d)|=|Z\cap(0,d)|+[g<d]-[d-1\in A],
\]
and $d-1\in A$ puts $d$ in $W$; so
\begin{equation}\label{eq:runsigns}
  w(d)=\varphi(d)|w(d)|\quad(d\le g),\qquad
  w(d)=-\varphi(d)|w(d)|\quad(d>g).
\end{equation}
In words: the $F_j$ keep the signs of $F$ at every integer, and $w$ keeps them
below $g$ and reverses them above.

\emph{The combination.}  At a run start $s_i$ we have $F_j(s_i)=0$ for $j\ne
i$, as $s_i\in Z_j$; $F_i(s_i)\ne0$, as $s_i\notin Z_i$; and $w(s_i)\ne0$,
as $s_i>g$ and $s_i-1\notin A$ keep $s_i$ out of $W$.  Put
$\rho_i=|w(s_i)|/|F_i(s_i)|$, $\mu=1/(1+\sum_i\rho_i)$, $\lambda_i=\mu\rho_i$
and
\[
  q=\sum_{j=1}^p\lambda_jF_j+\mu w .
\]
Then $q(0)=1$, and by \eqref{eq:runsigns}
$q(s_i)=\varphi(s_i)\bigl(\lambda_i|F_i(s_i)|-\mu|w(s_i)|\bigr)=0$.  Every
other point of $Z$ lies in $B_0$ or in $A\cap(A+1)$, hence in every $Z_j$ and
in $W$, so $q$ vanishes on all of $Z$.  Let $\omega$ be the weight of $w$ on
$y_{k+1}$.  It is nonzero, for otherwise $w$ would be a sum on $L-k-1$ nodes
with the $L-k-1$ zeros $W$; and since $y_{k+1}$ is the largest node and $w$ has
the sign $(-1)^{|W|}=(-1)^{L-k-1}$ past $\max W$, $\sgn\omega=(-1)^{L-k-1}$.
The sum $G=F-q$ on $y_{k+1},\ldots,y_L$ has weight $-\mu\omega\ne0$ on
$y_{k+1}$ and vanishes on the $L-k-1$ points $\{0\}\cup Z$, so by the zero
bound these are its only zeros and all are simple.  For large $d$ its sign is
$-\sgn\omega=(-1)^{L-k}=(-1)^{|Z|}$, which is $\varphi(d)$ there; as $G$ and
$\varphi$ change sign exactly at the points of $Z$,
\begin{equation}\label{eq:runG}
  \varphi(d)G(d)=|G(d)|\qquad(d\ge0).
\end{equation}

\emph{Decomposition.}  Put $\delta_d=|F(d)|-|q(d)|$ and
$\gamma_d=\sum_j\lambda_j|F_j(d)|+\mu|w(d)|-|q(d)|$, so that, all the series
converging absolutely, the left side of \eqref{eq:runchord} minus the right is
$D=\sum_{d\ge1}(\delta_d-\gamma_d)$.  For $1\le d\le g$ every term of $q(d)$
has the sign $\varphi(d)$ or vanishes, so $\gamma_d=0$ and
$\delta_d=\varphi(d)(F(d)-q(d))=\varphi(d)G(d)$.  For $d>g$ write
$P=\sum_j\lambda_j|F_j(d)|$ and $Q=\mu|w(d)|$; by \eqref{eq:runsigns}
$q(d)=\varphi(d)(P-Q)$, so $\gamma_d=P+Q-|P-Q|$ and
\[
  \delta_d-\gamma_d=|F(d)|-P-Q=\varphi(d)G(d)-2\mu|w(d)| ,
\]
whatever the sign of $P-Q$.  Hence
\[
  D=\sum_{d\ge1}\varphi(d)G(d)-2\mu\sum_{d>g}|w(d)| .
\]
Since $G(0)=0$, and $\varphi(d)\ne\varphi(d-1)$ only when $d-1\in Z$, where $G$
vanishes, $\sum_{d\ge1}\varphi(d)G(d)=\sum_{d\ge1}\varphi(d)G(d-1)$.  With
\eqref{eq:runsigns} and \eqref{eq:runG} this gives
\begin{equation}\label{eq:runsplit}
  \Tail(F)-\sum_j\lambda_j\Tail(F_j)-\mu\Tail(w)
  =\sum_{d=1}^{g}|G(d-1)|+\sum_{d>g}\varphi(d)E(d),
  \qquad E(d)=G(d-1)+2\mu\,w(d),
\end{equation}
and it remains to show $\varphi(d)E(d)\ge0$ for every integer $d>g$.

\emph{Zeros of $E$.}  $E$ is an exponential sum on $y_{k+1},\ldots,y_L$, and
its weight on $y_{k+1}$ is $\mu\omega(2-1/y_{k+1})$.  It vanishes at every $d$
with $d-1\in\{0\}\cup Z$ and $d\in W$.  With $S=\{0\}\cup B_0\cup\{g\}$ that
covers every $d$ with $d-1,d\in S$, and every $d\in A+1$.  Let $S$ consist of
the maximal runs $R_1,\ldots,R_m$ of consecutive integers, $0\in R_1$ and
$g\in R_m$.  At a gap start $q=r+1$ with $r\in S$, $r<g$, $q\notin S$, we have
$G(q-1)=0$ and $q\notin W$, so $E(q)=2\mu w(q)\ne0$, of sign $\varphi(q)$ by
\eqref{eq:runsigns}.

\emph{Counting.}  Suppose first $m\ge2$, with gap starts
$q_1<\cdots<q_{m-1}$, and $E\not\equiv0$.  The run $R_1$ supplies $|R_1|-1$
zeros below $q_1$; for
$1\le j\le m-2$ the run $R_{j+1}\subseteq B_0$ supplies $|R_{j+1}|-1$ zeros in
$(q_j,q_{j+1})$, and since
$\varphi(q_{j+1})=\varphi(q_j)(-1)^{|R_{j+1}|}$ that interval holds at least
$|R_{j+1}|$ zeros counted with multiplicity; and beyond $q_{m-1}$ lie the known
zeros $a+1,\ldots,g$, with $a=\min R_m$, and the $|A|$ points of $A+1$.  In
all $E$ has at least $|S|-2+|A|=|B_0|+|A|=L-k-2$ zeros, of which at least
$|R_m|-1+|A|$ lie in $(q_{m-1},\infty)$.  At $q_{m-1}$ the sign of $E$ is
$\varphi(q_{m-1})=(-1)^{|B_0|-|R_m|+1}$.  If $y_{k+1}<\tfrac12$ the weight on
$y_{k+1}$ is nonzero, $E$ has at most $L-k-1$ zeros, hence at most
$|R_m|+|A|$ in $(q_{m-1},\infty)$, and its sign at infinity is
$-\sgn\omega=(-1)^{|B_0|+|A|}$; the number of its zeros in
$(q_{m-1},\infty)$ is then congruent to $|R_m|-1+|A|$ modulo $2$, so it is
exactly $|R_m|-1+|A|$.  If $y_{k+1}=\tfrac12$ the weight vanishes, $E$ is a sum
on $L-k-1$ nodes with at most $L-k-2$ zeros, and again only the known zeros
occur.  Either way the zeros of $E$ beyond $q_{m-1}$ are exactly
$a+1,\ldots,g$ and the points of $A+1$, all simple.  So for an integer $d>g$
outside $A+1$, where $d-1\notin A$ and hence $|(A+1)\cap(0,d)|=|A\cap(0,d)|$,
\[
  \sgn E(d)=(-1)^{|B_0|-|R_m|+1}(-1)^{|R_m|-1}(-1)^{|A\cap(0,d)|}
  =(-1)^{|Z\cap(0,d)|}=\varphi(d),
\]
while $E$ vanishes on $A+1$.  If $m=1$, that is $B_0=\{1,\ldots,g-1\}$, the
known zeros $1,\ldots,g$ and $A+1$ number $L-k-1$.  For $y_{k+1}<\tfrac12$ they
are all the zeros, simple, and counting down from the sign
$(-1)^{|B_0|+|A|}$ at infinity gives $\sgn E(d)=\varphi(d)$ for $d>g$ outside
$A+1$ as before; for $y_{k+1}=\tfrac12$ the bound $L-k-2$ forces
$E\equiv0$.

So $\varphi(d)E(d)\ge0$ for every $d>g$, and \eqref{eq:runsplit} proves
\eqref{eq:runchord}.  Its first term vanishes exactly when $G$ vanishes on
$\{1,\ldots,g-1\}$, that is when $B_0=\{1,\ldots,g-1\}$; then $E\equiv0$ if
$y_{k+1}=\tfrac12$, and otherwise $\varphi(g+1)E(g+1)>0$, as $g+1\notin A+1$.
\end{proof}

Only the new node is held to the cutoff, and only through the sign of
$2-1/y_{k+1}$; the nodes $y_{k+2},\ldots,y_L$ may be any distinct reals in
$(0,y_{k+1})$; for $y_{k+1}>\tfrac12$ that sign flips.  The proof counts
integer slots: the zeros of $E$ pinned at consecutive pairs of
$\{0\}\cup B_0\cup\{g\}$, and its signs pinned at the gap starts, are what fix
its sign above $g$.  With the elements of $B_0$ allowed to be real the count
is unavailable, and \eqref{eq:runchord} is in fact false there, already for
$A=\{6\}$ in the confluent limit at $B_0=\{1,\tfrac52\}$, $g=3$.  Shifting
every run at once matters too: with a single neighbour the other runs of $A$
are not shifted in $w$, and $E$ loses the zeros they would force.

\begin{corollary}[The shift]\label{cor:shiftrun}
In the setting of Proposition~\ref{prop:runchord},
\begin{equation}\label{eq:shiftrun}
  \Tail\bigl(B_0\cup\{g\}\cup(A+1)\bigr)\le\Tail(B_0\cup A)
  \quad\text{whenever}\quad
  \Tail(B_0\cup A)\le\Tail\bigl(B_0\cup A\setminus\{s\}\cup\{e+1\}\bigr)
\end{equation}
for every run $[s,e]$ of $A$, the left side on $y_{k+1},\ldots,y_L$ and the
others on $y_{k+2},\ldots,y_L$.
\end{corollary}

\begin{proof}
By \eqref{eq:runchord} and the hypothesis,
$\Tail(F)\ge(1-\mu)\Tail(F)+\mu\Tail(w)$, and $\mu>0$.
\end{proof}

The hypothesis cannot simply be dropped.  In the confluent limit with all
nodes at $\tfrac12$, $B_0=\{1,\ldots,k\}$, $g=k+1$ and $A=\{M\}$, a
coin-tossing computation shows that the hypothesis and the conclusion of
\eqref{eq:shiftrun} are both equivalent to $M\ge2k+2$, the conclusion being
an equality exactly at $M=2k+2$.

The hypothesis is what minimality supplies.  For that we need
minimisers at every placement, and the argument of
Proposition~\ref{prop:attain} gives them.

\begin{lemma}[Full minimisers]\label{lem:attainX}
Let $0\le k<L$ and let $X$ be a set of at most $L-k-1$ positive integers.  Then
$\tau_k(X)$ is attained by a full certificate on $y_{k+1},\ldots,y_L$ whose
zero set contains $X$.
\end{lemma}

\begin{proof}
Run the first two steps of the proof of Proposition~\ref{prop:attain} on the
affine set $\{\lambda:\sum_i\lambda_i=1,\ v_x(\lambda)=0\ (x\in X)\}$ in place
of $\{\sum_i\lambda_i=1\}$.  It is closed, and nonempty because it contains the
full certificate of any zero set of size $L-k-1$ containing $X$, so the norm
$f$ attains its minimum on it.  If a minimiser has a zero set $Z\supseteq X$
with $|Z|\le L-k-2$, the conditions $\sum_ih_i=0$ and $v_d(h)=0$ for $d\in Z$
are at most $L-k-1$ linear conditions on $h\in\R^{L-k}$, so some $h\ne0$ meets
them, and the move $\lambda+th$ of that proof keeps the tail and every zero of
$Z$, and so every point of $X$, until it gains a zero.  The zero bound caps
the zeros at $L-k-1$.
\end{proof}

\begin{lemma}[Full minimisers at repeated roots]\label{lem:attainXrep}
Let $1<r_1\le\cdots\le r_L$, $0\le k<L$, and let $X$ be a set of at most
$L-k-1$ positive integers.  Then $\tau_k(X)$ is attained by a full
certificate on $\mathcal H_k$ whose zero set contains $X$.
\end{lemma}

\begin{proof}
The proof of Lemma~\ref{lem:attainX}, with the changes listed in the proof
of Proposition~\ref{prop:attainrep}.
\end{proof}

\begin{theorem}[The escaping placement is extremal]\label{thm:extremal}
Let $2\le r_1<\cdots<r_L$ and $0\le u\le L-1$.  Then \eqref{eq:step} holds for
every $0\le k<u$ and every set $X$ of $u-k$ positive integers.  Consequently
$T=\tau^*$, and
\[
  \inf_Qb_{u+1}\Bigl(Q\prod_{i=1}^L(x-r_i)\Bigr)=\frac1{\tau^*}
  =\min_M\max_{j\ge1}\bigl|S_j(P_uM)\bigr|
\]
over monic $Q,M\in\R[x]$.
\end{theorem}

\begin{proof}
Fix $k$ and $X$, and put $g=\max X$ and $X'=X\setminus\{g\}$.  By
Lemma~\ref{lem:attainX} on $y_{k+2},\ldots,y_L$, $\tau_{k+1}(X')$ is attained
by a full certificate $F$ with a zero set $Z^-\supseteq X'$ of size $L-k-2$.
If $g\in Z^-$ or $g>\max Z^-$, Proposition~\ref{prop:stepeasy} gives
\eqref{eq:step}.  Otherwise put $B_0=Z^-\cap(0,g)$, which contains $X'$, and
$A=Z^-\cap(g,\infty)\ne\varnothing$.  For a run $[s,e]$ of $A$, the set
$Z^-\setminus\{s\}\cup\{e+1\}$ has $L-k-2$ elements and contains $X'$, as
$s>g$, so its full certificate is a certificate for $X'$ and its tail is at
least $\tau_{k+1}(X')=\Tail(F)$.  The node $y_{k+1}=1/r_{k+1}$ is at most
$\tfrac12$, so Corollary~\ref{cor:shiftrun} applies: the full certificate $w$
on $y_{k+1},\ldots,y_L$ with zero set $B_0\cup\{g\}\cup(A+1)\supseteq X$ has
$\Tail(w)\le\Tail(F)$, whence $\tau_k(X)\le\Tail(w)\le\tau_{k+1}(X')$.
Proposition~\ref{prop:step} then gives $T=\tau^*$ and the value $1/\tau^*$,
and Proposition~\ref{prop:attain} makes \eqref{eq:multiplier} an equality at
some monic $M$.
\end{proof}

\paragraph{Repeated roots.}  For the rest of this subsection $2\le r_1\le\cdots\le r_L$, repetitions allowed, so every node is at most $1/2$.
Fix $0\le k<L-1$ and write $n=L-k-1=\dim\mathcal H_{k+1}$.  For
Lemma~\ref{lem:coalesce} applied to the nodes of $\mathcal H_k$, the largest
node is $y=1/r_{k+1}\le\tfrac12$, $\mathcal H$ is $\mathcal H_k$ and
$\mathcal H_-$ is $\mathcal H_{k+1}$; at level $\varepsilon$ the nodes
$N_\varepsilon$ are distinct, $y$ is the largest of them, and the others lie
in $(0,y)$, which is all that Theorem~2.2's deletion step and
Proposition~\ref{prop:runchord} ask of the nodes.  An object at level
$\varepsilon$ (a full certificate or a combination of such) carries a
superscript $\varepsilon$, and its limit, given by
Lemma~\ref{lem:coalesce}, the superscript $0$.

\begin{proposition}[The step in the aligned cases, repeated roots]
\label{prop:stepeasyrep}
Let $X$ be finite and nonempty, $g=\max X$, $X'=X\setminus\{g\}$, and let
$v^-$ be a full certificate on $\mathcal H_{k+1}$ whose zero set $Z^-$
contains $X'$.  If $g\in Z^-$ or $g>\max Z^-$, then some certificate for $X$
on $\mathcal H_k$ has tail at most $\Tail(v^-)$.
\end{proposition}

\begin{proof}
Form $Z^+$ as in Proposition~\ref{prop:stepeasy}; its largest element
exceeds the others and $Z^+\setminus\{\max Z^+\}=Z^-$.  At level
$\varepsilon$ the full certificates $v^{+,\varepsilon}$ on $N_\varepsilon$
with zero set $Z^+$ and $v^{-,\varepsilon}$ on $N^-_\varepsilon$ with zero
set $Z^-$ satisfy $\Tail(v^{+,\varepsilon})\le\Tail(v^{-,\varepsilon})$: by
the deletion step of Theorem~2.2 if $Z^+$ is not $\{1,\ldots,n\}$, the
largest node being at most $\tfrac12$, and by Lemma~2.1, which gives the
ratio $y/(1-y)\le1$, if it is.  By Lemma~\ref{lem:coalesce} both tails
converge, to $\Tail(v^+)$ and $\Tail(v^-)$, where $v^+$ is the full
certificate on $\mathcal H_k$ with zero set $Z^+\supseteq X$.
\end{proof}

\begin{proposition}[The run chord at repeated roots]\label{prop:runchordrep}
Let $g\ge1$, $B_0\subseteq\{1,\ldots,g-1\}$ and $A$ a nonempty set of
integers above $g$ with $|B_0|+|A|=n-1$, with runs $[s_j,e_j]$ and
neighbours $Z_j$ as before Proposition~\ref{prop:runchord}.  Let $F$ and
$F_j$ be the full certificates on $\mathcal H_{k+1}$ with zero sets
$Z=B_0\cup A$ and $Z_j$, and $w$ the full certificate on $\mathcal H_k$ with
zero set $W=B_0\cup\{g\}\cup(A+1)$.  Then there are
$\lambda_1,\ldots,\lambda_p,\mu>0$ with $\sum_j\lambda_j+\mu=1$ and
\[
  \Tail(F)\ \ge\ \sum_{j=1}^p\lambda_j\Tail(F_j)+\mu\,\Tail(w),
\]
with equality if and only if $r_{k+1}=2$ and $B_0=\{1,\ldots,g-1\}$.
\end{proposition}

\begin{proof}
At level $\varepsilon$ Proposition~\ref{prop:runchord} applies, with the
weights of its proof: $\rho^\varepsilon_i=|w^\varepsilon(s_i)|/
|F_i^\varepsilon(s_i)|$, $\mu^\varepsilon=1/(1+\sum_i\rho^\varepsilon_i)$,
$\lambda^\varepsilon_i=\mu^\varepsilon\rho^\varepsilon_i$.  By
Lemma~\ref{lem:coalesce}, $w^\varepsilon(s_i)\to w^0(s_i)$ and
$F_i^\varepsilon(s_i)\to F_i^0(s_i)$, where $w^0=w$ and $F_i^0=F_i$ are the
full certificates above.  Both limits are nonzero by the zero bound, since
$s_i\notin W$ and $s_i\notin Z_i$ (see the proof of
Proposition~\ref{prop:runchord}).  So the weights converge to
$\lambda_i=\mu\rho_i>0$ and $\mu=1/(1+\sum_i\rho_i)>0$, and the tails
converge; this proves the inequality.

For the equality clause pass to the limit in \eqref{eq:runsplit}.  With
$q^\varepsilon=\sum_j\lambda^\varepsilon_jF^\varepsilon_j+\mu^\varepsilon
w^\varepsilon$, $G^\varepsilon=F^\varepsilon-q^\varepsilon$ and
$E^\varepsilon(d)=G^\varepsilon(d-1)+2\mu^\varepsilon w^\varepsilon(d)$,
that identity and the proof of Proposition~\ref{prop:runchord} give
\[
  \Tail(F^\varepsilon)-\sum_j\lambda^\varepsilon_j\Tail(F^\varepsilon_j)
  -\mu^\varepsilon\Tail(w^\varepsilon)
  =\sum_{d=1}^g|G^\varepsilon(d-1)|+\sum_{d>g}\varphi(d)E^\varepsilon(d),
  \qquad\varphi(d)E^\varepsilon(d)\ge0 .
\]
All terms converge, the last sum by dominated convergence
(Lemma~\ref{lem:coalesce}), so the defect $\Delta$ of the proposition equals
$\sum_{d=1}^g|G^0(d-1)|+\sum_{d>g}\varphi(d)E^0(d)$ with every
$\varphi(d)E^0(d)\ge0$.  Let $\omega$ be the coordinate of $w$ on the top
basis function $\psi_{k+1}$ of $\mathcal H_k$; it is nonzero, since otherwise
$w\in\mathcal H_{k+1}$, of dimension $n$, would have the $n$ zeros $W$.  The
$F$'s lie in $\mathcal H_{k+1}$, so $G^0\in\mathcal H_k$ has coordinate
$-\mu\omega\ne0$ on $\psi_{k+1}$; it vanishes on the $n$ points
$\{0\}\cup Z$ and so nowhere else.  Hence the first sum is positive unless
$\{1,\ldots,g-1\}\subseteq Z$, that is $B_0=\{1,\ldots,g-1\}$.  In that
case, if $r_{k+1}>2$: the shift $d\mapsto G^0(d-1)$ maps $\mathcal H_k$ to
itself and multiplies the coordinate on $\psi_{k+1}=d^{e-1}y^d$ by $1/y$,
because $(d-1)^ay^{d-1}=y^{-1}(d^a+\text{lower powers})y^d$; so $E^0$ has
coordinate $\mu\omega(2-1/y)\ne0$ on $\psi_{k+1}$ and at most $n$ zeros.  It
vanishes at $1,\ldots,g$ and on $A+1$, which are $g+|A|=n$ points, so
$E^0(g+1)\ne0$, as $g+1\notin A+1$, and $\varphi(g+1)E^0(g+1)>0$.  If
$r_{k+1}=2$ and $B_0=\{1,\ldots,g-1\}$, every level is an equality by
Proposition~\ref{prop:runchord}, and so is the limit.
\end{proof}

\begin{corollary}[The shift at repeated roots]\label{cor:shiftrep}
In the setting of Proposition~\ref{prop:runchordrep}, if
$\Tail(F)\le\Tail(F_j)$ for every run, then $\Tail(w)\le\Tail(F)$.
\end{corollary}

\begin{proof}
As for Corollary~\ref{cor:shiftrun}:
$\Tail(F)\ge(1-\mu)\Tail(F)+\mu\Tail(w)$ with $\mu>0$.
\end{proof}

\begin{theorem}[The escaping placement is extremal at repeated roots]
\label{thm:extremalrep}
Let $2\le r_1\le\cdots\le r_L$, repetitions allowed, and $0\le u\le L-1$.
Then \eqref{eq:step} holds on the Hermite spaces, that is
$\tau_k(X)\le\tau_{k+1}(X\setminus\{\max X\})$ for every $0\le k<u$ and every
set $X$ of $u-k$ positive integers.  Consequently $T=\tau^*$, and
\[
  \inf_Qb_{u+1}\Bigl(Q\prod_{i=1}^L(x-r_i)\Bigr)=\frac1{\tau^*}
  =\min_M\max_{j\ge1}\bigl|S_j(P_uM)\bigr|
\]
over monic $Q,M\in\R[x]$, with $\tau^*$ the least tail of a certificate on
the Hermite space $\mathcal H_u$ of the surviving roots.
\end{theorem}

\begin{proof}
Word for word the proof of Theorem~\ref{thm:extremal}, with
Lemma~\ref{lem:attainXrep}, Proposition~\ref{prop:stepeasyrep} and
Corollary~\ref{cor:shiftrep} in place of Lemma~\ref{lem:attainX},
Proposition~\ref{prop:stepeasy} and Corollary~\ref{cor:shiftrun}: the
minimiser $F$ for $\tau_{k+1}(X')$ is full on $\mathcal H_{k+1}$ with zero set
$Z^-\supseteq X'$; in the aligned cases Proposition~\ref{prop:stepeasyrep}
gives \eqref{eq:step}; otherwise each neighbour $Z^-\setminus\{s\}\cup\{e+1\}$
contains $X'$, so minimality gives the hypothesis of
Corollary~\ref{cor:shiftrep}, and $w$ is a certificate for $X$ on
$\mathcal H_k$ with $\Tail(w)\le\Tail(F)$.  Chaining as in
Proposition~\ref{prop:step} gives $\tau(X)\le\tau^*$ for all $|X|=u$, so
$T=\tau^*$ by Theorem~\ref{thm:valuerep}, and
Proposition~\ref{prop:attainrep} gives the multiplier form.
\end{proof}

In particular the proof of Corollary~\ref{cor:charrep} may cite
Theorem~\ref{thm:valuerep} for its first sentence, and when
\eqref{eq:partial} fails the infimum at repeated roots at least $2$ is
$1/\tau^*$.

So when \eqref{eq:partial} fails the infimum is always $1/\tau^*$, and
Corollary~\ref{cor:uone} reads $\inf_Qb_2=1/\tau^*$.  The minimizing zero
set need not be a consecutive block with one displaced zero.  At
$(2,3,4,5,6)$ with $u=0$ it is $\{1,2,4,6\}$, of tail $\tfrac{41}{12600}$,
against $0.003847$ for the best block $\{1,2,3\}$ with one displaced zero.

\paragraph{Naive steps fail.}  The shift of every run cannot be replaced by a
simpler move.  Replacing $\max Z^-$ by $g$ throws away the displacement
$Z^-$ earned and can leave a consecutive set, where Theorem~2.2 is an
equality: at $(\tfrac{503}{250},\tfrac{2053}{1000},\tfrac{527}{250},
\tfrac{11}5,\tfrac{1177}{500})$ with $u=1$ and $X=\{2\}$ that turns
$Z^-=\{1,3,6\}$ into $\{1,2,3\}$, and no displaced lift of it comes within a
factor $1.40$.  Adjoining $g$ to $Z^-$ and displacing nothing gives a zero
set of the right size, but at $(\tfrac{2013}{1000},\tfrac{257}{125},
\tfrac{2127}{1000},\tfrac{1071}{500},\tfrac{549}{250})$ with $u=1$ and
$X=\{2\}$ its tail exceeds that of $Z^-$ by a factor $2.01$.  Sliding the
smallest zero $M$ of $Z^-$ above $g$ down to $g$, and letting one further zero
move freely, fails as well, even when $M$ is the best choice with
$B=Z^-\setminus\{M\}$ fixed: at $L=9$, $k=0$,
$y_i=\tfrac12-(i-1)\cdot10^{-4}$, $B=\{1,2,3,11,12,13\}$, $g=4$ and $M=6$,
every certificate on all nine nodes vanishing on $B\cup\{4\}$ has tail at
least $1.026$ times that of the full certificate on $y_2,\ldots,y_9$ with zero
set $B\cup\{6\}$, and in the confluent limit with all nine nodes at
$\tfrac12$, where the full certificate with zero set $Z$ is
$d\mapsto2^{-d}\prod_{z\in Z}(1-d/z)$, the ratio is exactly
$\tfrac3{143}\big/\tfrac{35}{1716}=\tfrac{36}{35}$ (exact rational
arithmetic).  There $B\cup\{6\}$ is not a minimizer: the minimizer for
$X'=\{1,2,3\}$ on $y_2,\ldots,y_9$ is $\{1,2,3,6,8,12,17\}$, and
\eqref{eq:step} holds at $X=\{1,2,3,4\}$ with ratio $0.70$.  Minimality of
$Z^-$ is what the proof of Theorem~\ref{thm:extremal} uses.

When $Z^-\cap(0,g)$ is an initial segment $\{1,\ldots,k'\}$ an explicit
certificate does the step, without minimality.  Let $F(x)=\sum_dF_dx^d$ be
the generating function of any full certificate with such a zero set
$Z^-\not\ni g$, and $y\le\tfrac12$ the new node.  Then
$v(x)=(1-yxF(x))/(1-yx)$ is an exponential sum on all the nodes with
$v_0=1$ and $v_d=-\sum_{1\le j<d}F_jy^{d-j}$, so it vanishes on
$\{1,\ldots,k'+1\}$ and $\Tail(v)\le\tfrac{y}{1-y}\Tail(F)\le\Tail(F)$.
For $g>k'+1$, $F$ keeps one sign on $(k',g]$, so $v_g$ and $F_g$ have
opposite signs, and the segment from $F$ to $v$ crosses $v_g=0$ at a
certificate vanishing on $\{1,\ldots,k'\}\cup\{g\}$ with tail at most
$\Tail(F)$, by convexity.  Only the new node is held to the cutoff.

\subsection*{Two examples}

For $u\ge2$ the covering of Theorem~\ref{thm:lift} is also a finite
certificate of the value.  Start from $(A,k)=(\varnothing,u)$.  Choose a full
certificate on $y_{k+1},\ldots,y_L$ whose zero set $Z$ contains $A$ and whose
tail is at most $\tau^*$.  By Theorem~\ref{thm:lift} it covers every
$X\supseteq A$ whose other $k$ elements lie in $Z\cup(\max Z,\infty)$.
Recurse on $(A\cup\{w\},k-1)$ for each $w<\max Z$ outside $Z\cup A$.  By
Theorem~\ref{thm:extremal} the search always closes, and where the lifts are
strict it also shows that the infimum is not attained.

\begin{proposition}[Two infima]\label{prop:values}
Over monic $Q\in\R[x]$,
\begin{align*}
  \inf_Qb_2\bigl(Q\,(x-2)(x-3)(x-5)(x-7)\bigr)
    &=\frac{1704}{31}=54.967\ldots,\\
  \inf_Qb_3\Bigl(Q\prod_{r=2}^7(x-r)\Bigr)
    &=\frac{114840}{259}=443.397\ldots,
\end{align*}
against the bounds $48$ and $360$ of \eqref{eq:order} and the values $57$ and
$480$ of the construction.  Neither infimum is attained.
\end{proposition}

\begin{proof}
At $(2,3,5,7)$ the full certificate on $(\tfrac13,\tfrac15,\tfrac17)$ with
zero set $\{1,3\}$ is
$v^*_d=(81\cdot3^{-d}-625\cdot5^{-d}+686\cdot7^{-d})/142$.  It has
$v^*_2=-\tfrac1{71}$, is positive past $3$, and has tail $\tfrac{31}{1704}$.
The multiplier $M=x+\tfrac9{62}$ gives
\[
  \frac{\widetilde P_1(y)\,(1+\tfrac9{62}y)}{1-y}
  =1-\tfrac{859}{62}y+\tfrac{1704}{31}y^2-\tfrac{2463}{62}y^3
   -\tfrac{1704}{31}\sum_{d\ge4}y^d ,
\]
so by \eqref{eq:dualpair} $\tau^*=\tfrac{31}{1704}$, with $Z^*=\{1,3\}$ and
$G=\{2\}$.  The full certificate on all four nodes with zero set $\{1,2,5\}$
has tail $\tfrac{21337}{1910352}<\tfrac{31}{1704}$, and
Corollary~\ref{cor:uone} gives the first value.

At $(2,\ldots,7)$ with $u=2$ the full certificate on
$(\tfrac14,\ldots,\tfrac17)$ with zero set $\{1,2,4\}$ has
$v^*_3=\tfrac1{638}$, is negative past $4$, and has tail
$\tfrac{259}{114840}$.  The multiplier $M=x+\tfrac{60}{259}$ gives the
partial sums $1$, $-\tfrac{5379}{259}$, $\tfrac{5666}{37}$,
$-\tfrac{114840}{259}$, $\tfrac{64440}{259}$ and then $\tfrac{114840}{259}$
forever, so $\tau^*=\tfrac{259}{114840}$ with $Z^*=\{1,2,4\}$.  The search
closes in two steps.  Theorem~\ref{thm:lift} with $k=2$ covers every
$X\not\ni3$.  The full certificate on $(\tfrac13,\ldots,\tfrac17)$ with zero
set $\{1,2,3,5\}$, of tail $\tfrac{31}{36720}$, covers every $\{3,x\}$ except
$\{3,4\}$.  The full certificate on all six nodes with zero set
$\{1,2,3,4,7\}$, of tail $\tfrac{119261}{193623120}$, covers $\{3,4\}$.

Neither infimum is attained, because $\tau(X)<\tau^*$ for every $X$.  Each
covering above either lifts with $k\ge1$ onto a padded zero set that is not
consecutive, which is strict by Theorem~\ref{thm:lift}, or has tail below
$\tau^*$ outright.  So every multiple $F$ has
$b_{u+1}(F)\ge1/\tau(X_F)>1/\tau^*$.
\end{proof}

\begin{proposition}[An infimum at a repeated root]\label{prop:valuesrep}
Over monic $Q\in\R[x]$,
\[
  \inf_Qb_2\bigl(Q\,(x-2)(x-3)^2(x-5)\bigr)=24,
  \qquad
  \inf_Qb_2\bigl(Q\,(x-2)(x-3)^3\bigr)=\tfrac{27}2,
\]
against the bounds $16$ and $8$ of \eqref{eq:order} and the values $29$ and
$19$ of the construction ($\max_j|S_j(P_1)|$).
\end{proposition}

\begin{proof}
Here $u=1$ and $\mathcal H_1$ is the Hermite space of $(3,3,5)$, resp.\
$(3,3,3)$.  The full certificate with zero set $\{1,3\}$ is
\[
  v^*_d=\frac{(24d-99)3^{-d}+125\cdot5^{-d}}{26},
  \qquad\text{resp.}\qquad
  v^*_d=\frac{(d-1)(d-3)}3\,3^{-d};
\]
the first has $v^*_2=-\tfrac1{39}$ and is positive past $3$, the second is
$-\tfrac1{27}$ at $2$ and positive past $3$, and summing
$\sum_dd^ay^d$ in closed form gives the tails $\tfrac1{24}$ and
$\tfrac2{27}$.  The multipliers $M=x+\tfrac12$ and $M=x+\tfrac{11}{16}$
give the partial sums
\[
  1,-\tfrac{19}2,24,-\tfrac32,-24,-24,\ldots
  \qquad\text{and}\qquad
  1,-\tfrac{117}{16},\tfrac{27}2,\tfrac{81}{16},-\tfrac{27}2,-\tfrac{27}2,\ldots,
\]
so by \eqref{eq:dualpair} on $\mathcal H_1$, $\tau^*=\tfrac1{24}$ and
$\tfrac2{27}$, and Theorem~\ref{thm:extremalrep} gives the values.
\end{proof}

\paragraph{What the limit does not give.}  Theorem~\ref{thm:extremalrep}
needs only non-strict inequalities, which is why the coalescence limit
suffices.  Whether these infima are attained would need the strict clause of
Theorem~\ref{thm:lift} at repeated roots, which we have not carried over; an
exact search finds, at both root sets, a certificate with tail below
$\tau^*$ for every placement $X\subseteq\{1,\ldots,30\}$.

In both cases the extremal series is the construction of
Proposition~\ref{prop:partial} applied to $P_uM$ rather than $P_u$, and the
factor $x+\beta$ has a negative root.  It raises the total $|P_u(1)M(1)|$
above $|P_u(1)|$ but brings every partial sum within it, so \eqref{eq:partial}
holds for the product.  Complementary slackness explains the shape.  Past the
largest zero $g$ of $v^*$ the witness is constant, so $(1-y)W$ is a
polynomial of degree $g+1$ and $\deg M=g+1-m$, which is $1$ both times.

\section{Below the cutoff}\label{sec:subtwo}

This section decides when the bounds \eqref{eq:repaired} and
\eqref{eq:twosided} of Section~\ref{sec:extend} are the infimum.  They
differ from \eqref{eq:order} only when an exempt root lies below $2$.

At $(\tfrac32,3)$ Corollary~\ref{cor:repaired} gives $b_2\ge1$, below the $1.62$ of
Proposition~\ref{prop:subtwo}, and $1$ is the infimum.  More generally the
repaired bound is attained whenever the $u$ exempt roots lie in $(1,2]$ and
the partial sums of the whole root product are small enough.  The multiples
that attain it put the exempt coefficients at the top rather than the
bottom.  Here $S_j(P)$ are the partial sums of \eqref{eq:wgen} with $u=0$.

\begin{proposition}[Sharpness below $2$]\label{prop:subtwosharp}
Let $r_1\le\cdots\le r_L$ be real with $r_1>1$, repetitions allowed, put
$P=\prod_{i=1}^L(x-r_i)$, and let $0\le u\le L-1$.  Let $\sigma$ be the
$(u+1)$-st largest of $|S_1(P)|,\ldots,|S_{L-1}(P)|$, or $0$ if $u=L-1$.  Then
\[
  \inf_Qb_{u+1}(QP)\le\max\bigl(|P(1)|,\sigma\bigr)
\]
over monic $Q\in\R[x]$.  If $r_u\le2$ and $\sigma\le|P(1)|$, the infimum is
$|P(1)|$, which is the value of \eqref{eq:repaired}.  In particular
$\inf_Qb_L(QP)=|P(1)|$ whenever $r_{L-1}\le2$.
\end{proposition}

\begin{proof}
Since $P(0)\ne0$, $x$ is invertible modulo $P$.  For $n\ge1$ let $R_n$ be the
polynomial of degree less than $L$ with $x^nR_n\equiv1\pmod P$.  Its
coefficient vector is $T^n$ applied to that of $1$, where $T$ is
multiplication by $x^{-1}$ on $\R[x]/(P)$.  The eigenvalues of $T$ are the
$1/r_i<1$, so $R_n\to0$, and for large $n$ we have $R_n(1)\ne1$.  Put
\[
  t_n=\frac{P(1)}{1-R_n(1)},\qquad M_n=P+t_nR_n .
\]
Then $M_n$ is monic of degree $L$ with $M_n(1)=t_n$, so
$M_n=(x-1)\widetilde A_n+t_n$ with $\widetilde A_n$ monic of degree $L-1$.
Moreover $x^nM_n-t_n\equiv x^nP+t_n(x^nR_n-1)\equiv0\pmod P$, and
$x^nM_n-t_n$ vanishes at $1$, where $P$ does not, so $P$ divides
\[
  F_n=\frac{x^nM_n-t_n}{x-1}=x^n\widetilde A_n+t_n\sum_{j=0}^{n-1}x^j .
\]
Thus $F_n$ is a monic multiple of $P$.  Its nonleading coefficients are the
$L-1$ lower coefficients of $\widetilde A_n$, followed by $t_n$ repeated $n$
times.  At most $u$ of those can exceed the larger of $|t_n|$ and the
$(u+1)$-st largest lower coefficient of $\widetilde A_n$, so $b_{u+1}(F_n)$ is
at most that larger value.  As $n\to\infty$, $t_n\to P(1)$ and
$\widetilde A_n\to(P-P(1))/(x-1)$.  By \eqref{eq:wgen} the lower coefficients
of that limit are $S_1(P),\ldots,S_{L-1}(P)$, which proves the bound.  If
$r_u\le2$, every $\min(r_i-1,1)$ with $i\le u$ is $r_i-1$, so the right side
of \eqref{eq:repaired} is $|P(1)|$.  At $u=L-1$, $\sigma=0$.
\end{proof}

At $(\tfrac32,10,20)$ with $u=1$, for example, the partial sums are
$1,-30.5,214.5,-85.5$, so $\sigma=30.5$ and the infimum of $b_2$ is
$85.5=|P(1)|$, while \eqref{eq:order} would assert $171$.

At $(\tfrac32,\tfrac52,4)$ with $u=2$, \eqref{eq:repaired} gives $b_3\ge\tfrac32$
and \eqref{eq:twosided} gives $b_3\ge\tfrac94=|P(1)|$, which is the infimum by
Proposition~\ref{prop:subtwosharp} ($\sigma=0$ at $u=L-1$).  Whether
\eqref{eq:twosided} is attained is decided by the two constructions already in
hand, one for each side of the minimum.

\begin{theorem}[When the two-sided bound is the infimum]\label{thm:twosidedsharp}
Let $1<r_1<\cdots<r_L$ be distinct and $0\le u\le L-1$, put
$\pi=\prod_{i\le u}(r_i-1)$, and let $\sigma$ be as in
Proposition~\ref{prop:subtwosharp}.  Then, over monic $Q\in\R[x]$,
\[
  \inf_Qb_{u+1}(QP)=\min\bigl(|P(1)|,|P_u(1)|\bigr)
\]
if and only if either $\pi\le1$ and $\sigma\le|P(1)|$, or $\pi\ge1$ and
\eqref{eq:partial} holds for $P_u$.  Otherwise the infimum is strictly larger.
\end{theorem}

\begin{proof}
Write $R=\min(|P(1)|,|P_u(1)|)$; Corollary~\ref{cor:twosided} gives
$\inf\ge R$.  If $\pi\le1$ then $R=|P(1)|$, and if also $\sigma\le|P(1)|$,
Proposition~\ref{prop:subtwosharp} gives $\inf\le|P(1)|$.  If $\pi\ge1$ then
$R=|P_u(1)|$, and if \eqref{eq:partial} holds, \eqref{eq:multiplier} at $M=1$
gives $\inf\le\max_{j\ge1}|S_j(P_u)|=|P_u(1)|$, since the partial sums
stabilize at $S_m(P_u)=P_u(1)$.  Neither step uses $r_1\ge2$:
Theorem~\ref{thm:value} and \eqref{eq:multiplier} assume only $1<r_1$.

Conversely suppose neither condition holds.  By \eqref{eq:weak} every monic
multiple has $b_{u+1}\ge1/\tau(X_F)$, so it suffices to show
$T=\sup_{|X|=u}\tau(X)<1/R$.  Exactly as in the proof of
Theorem~\ref{thm:converse}, a sequence $X_\nu$ with $\tau(X_\nu)\to1/R$ has a
subsequence $X_\nu=X_0\cup E_\nu$ with $|E_\nu|=k$ and $\min E_\nu\to\infty$,
and Proposition~\ref{prop:farrep}(c), which needs only $r_1>1$, turns a
certificate for $X_0$ on $y_{k+1},\ldots,y_L$ with tail below $1/R$ into
certificates for $X_\nu$ whose tails tend to it.  So it suffices to find, for
every $0\le k\le u$ and every $X_0$ of size $u-k$, such a certificate.

Take the zero set $Z$ of $X_0$ and the first $m-1$ positive integers outside
it, of size $L-k-1$ and leading run $f\ge m-1$.  On the nodes
$y_{k+1},\ldots,y_L$ the products $\Pi_j$ are the same as on all $L$ nodes for
$j\le L-k$, so Theorem~\ref{thm:twosided} bounds the tail by
$\max(\Pi_{f+1},\Pi_{L-k})$, strictly unless $Z$ is consecutive.  With
distinct roots the increments of $\log\Pi_j$ increase strictly, so
$j\mapsto\log\Pi_j$ is strictly convex, and since $m\le f+1\le L-k\le L$,
\[
  \max(\Pi_{f+1},\Pi_{L-k})\le\max(\Pi_m,\Pi_L)=1/R,
\]
with $\Pi_{L-k}<1/R$ unless $L-k$ is an endpoint $m$ or $L$ carrying the
maximum.  So the tail is below $1/R$ except when $Z=\{1,\ldots,L-k-1\}$ and
either $k=u$ and $\Pi_m\ge\Pi_L$, or $k=0$ and $\Pi_L\ge\Pi_m$.

If $k=u$ and $\Pi_m\ge\Pi_L$, that is $\pi\ge1$, then $X_0=\varnothing$ and the
tail is $1/|P_u(1)|=1/R$; as $\pi\ge1$, \eqref{eq:partial} fails for $P_u$, and
Lemma~\ref{lem:improve} with $Q=P_u$ and $X_0=\varnothing$ gives a certificate
with tail below $1/R$.  If $k=0$ and $\Pi_L\ge\Pi_m$, that is $\pi\le1$, then
$X_0\subseteq\{1,\ldots,L-1\}$ and the tail is $1/|P(1)|=1/R$; as $\pi\le1$,
$\sigma>|P(1)|$, so at least $u+1$ indices $j\in\{1,\ldots,L-1\}$ have
$|S_j(P)|>|P(1)|$, one of them outside $X_0$, and Lemma~\ref{lem:improve} with
$Q=P$ gives a certificate for $X_0$ with tail below $1/R$.  (For $u=0$ both
cases are $k=0$, $X_0=\varnothing$, and either argument applies.)
\end{proof}

For $r_1\ge2$ one has $\pi\ge1$, with $\pi=1$ only at $u=0$ or at $u=1$,
$r_1=2$, and the theorem is Proposition~\ref{prop:partial} with
Theorem~\ref{thm:converse}: at $u=1$, $r_1=2$ the proof of
Theorem~\ref{thm:converse} shows that $\sigma\le|P(1)|$ already forces
\eqref{eq:partial}.  Below $2$ it answers the question left by
Proposition~\ref{prop:subtwosharp}.

\begin{corollary}[When the repaired bound is the infimum]\label{cor:repairedsharp}
Let $1<r_1<\cdots<r_L$ be distinct and $1\le u\le L-1$.  Then
\eqref{eq:repaired} is the infimum of $b_{u+1}$ over monic multiples if and
only if either $r_u\le2$ and $\sigma\le|P(1)|$, or $r_1\ge2$ and
\eqref{eq:partial} holds.  In particular, with an exempt root below $2$ it is
the infimum exactly under the hypotheses of
Proposition~\ref{prop:subtwosharp}; when $r_1<2<r_u$ it is never the
infimum, even with repeated roots.
\end{corollary}

\begin{proof}
If $r_1<2<r_u$, Corollary~\ref{cor:twosided} puts every $b_{u+1}$ at or above
a value strictly larger than \eqref{eq:repaired}; that argument allows
repeated roots.  If $r_u\le2$ and $r_1<2$, then \eqref{eq:repaired} is
$|P(1)|=\min(|P(1)|,|P_u(1)|)$ and $\pi\le r_1-1<1$, so by
Theorem~\ref{thm:twosidedsharp} it is the infimum exactly when
$\sigma\le|P(1)|$.  If $r_1\ge2$ it is \eqref{eq:order}, settled by
Proposition~\ref{prop:partial} and Theorem~\ref{thm:converse}.
\end{proof}

\begin{theorem}[When the two-sided bound is the infimum, repeated roots]
\label{thm:twosidedsharprep}
Theorem~\ref{thm:twosidedsharp} holds verbatim for $1<r_1\le\cdots\le r_L$,
repetitions allowed: with $\pi=\prod_{i\le u}(r_i-1)$ and $\sigma$ as in
Proposition~\ref{prop:subtwosharp},
$\inf_Qb_{u+1}(QP)=\min(|P(1)|,|P_u(1)|)$ if and only if either $\pi\le1$
and $\sigma\le|P(1)|$, or $\pi\ge1$ and \eqref{eq:partial} holds for $P_u$;
otherwise the infimum is strictly larger.
\end{theorem}

\begin{proof}
Write $R=\min(|P(1)|,|P_u(1)|)$; Corollary~\ref{cor:twosided} gives
$\inf\ge R$ at repeated roots.  Sufficiency: if $\pi\le1$ and
$\sigma\le|P(1)|$, Proposition~\ref{prop:subtwosharp}, which allows
repetitions, gives $\inf\le|P(1)|=R$; if $\pi\ge1$ and \eqref{eq:partial}
holds, Theorem~\ref{thm:valuerep} and Proposition~\ref{prop:attainrep} with
$M=1$ give $\inf\le1/\tau^*\le\max_j|S_j(P_u)|=|P_u(1)|=R$.

Necessity.  Suppose neither condition holds.  As in the proof of
Theorem~\ref{thm:twosidedsharp}, by \eqref{eq:weak} on $\mathcal H_0$ and
Proposition~\ref{prop:farrep}(c), which allows repetitions and needs only
$r_1>1$, it suffices to find, for every $0\le k\le u$ and every $X_0$ of size
$u-k$, a certificate for $X_0$ on $\mathcal H_k$ with tail below $1/R$.  Let
$Z$ consist of $X_0$ and the first $m-1$ positive integers outside it,
$m=L-u$; it has $L-k-1$ elements and leading run $f\ge m-1$.  In increasing
order the nodes of $\mathcal H_k$ are $1/r_L\le\cdots\le1/r_{k+1}$, so the
products $\Pi_j=\prod_{i>L-j}(r_i-1)^{-1}$, $j\le L-k$, are those of all $L$
roots, and Proposition~\ref{prop:twosidedconf} with $t=0$ bounds the tail of
the full certificate by $\max(\Pi_{f+1},\Pi_{L-k})$, strictly unless
$Z=\{1,\ldots,L-k-1\}$.  The increments
$\delta_j=\log\Pi_j-\log\Pi_{j-1}=-\log(r_{L-j+1}-1)$ are nondecreasing in
$j$, so $\log\Pi_j$ is convex and
$\max(\Pi_{f+1},\Pi_{L-k})\le\max(\Pi_m,\Pi_L)=1/R$, since
$m\le f+1\le L-k\le L$.  So the tail is below $1/R$ unless $Z$ is
consecutive and $\Pi_{L-k}=\max(\Pi_m,\Pi_L)$.

\emph{Endpoints.}  If $k=u$ and $\Pi_m\ge\Pi_L$, or $k=0$ and
$\Pi_L\ge\Pi_m$, the last two paragraphs of the proof of
Theorem~\ref{thm:twosidedsharp} apply unchanged, with
Lemma~\ref{lem:improve} with repetitions.

\emph{Interior.}  Let $0<k<u$ with $\Pi_{L-k}=\max(\Pi_m,\Pi_L)$, and put
$j=L-k$.  If $\delta_{j+1}>0$ then $\delta_i>0$ for all $i>j$ and
$\Pi_L>\Pi_j$; so $\delta_{j+1}\le0$, hence $\delta_i\le0$ for
$m<i\le j$, and $\Pi_j\ge\Pi_m$ forces $\delta_i=0$ there; then
$\delta_{j+1}\ge\delta_j=0$, so $\delta_{j+1}=0$, and $\Pi_L\le\Pi_j$ with
nonnegative $\delta_i$ for $i>j$ forces those to vanish too.  So every
$\delta_i$, $m<i\le L$, is $0$: $r_1=\cdots=r_u=2$.  Then $\pi=1$, and as the
conditions fail, \eqref{eq:partial} fails for $P_u$, all roots are at least
$2$, and $Z=\{1,\ldots,L-k-1\}$ contains $X_0$.  This is the equality case of
the proof of Theorem~\ref{thm:converse}: with
$Q=\prod_{i>k}(x-r_i)=(x-2)^{u-k}P_u$, the set $\mathcal A(Q)$ there contains
$u-k+1$ integers of $\{1,\ldots,L-k-1\}$, one of them outside $X_0$, and
Lemma~\ref{lem:improve} with repetitions gives a certificate for $X_0$ on
$\mathcal H_k$ with tail below $1/|Q(1)|=1/|P_u(1)|=1/R$.
\end{proof}

\paragraph{The tie.}  The interior case is empty for distinct roots, where
the $\delta_j$ increase strictly, and the proof of
Theorem~\ref{thm:twosidedsharp} asserts $\Pi_{L-k}<1/R$ for $0<k<u$.  At
repeated roots that assertion fails, but only at $r_1=\cdots=r_u=2$, where
the criterion itself does not change.  At $(2,2,3,5,7)$ with $u=2$, $k=1$ and
$X_0=\{1\}$, the full certificate on the Hermite space of $(2,3,5,7)$ with
the consecutive zero set $\{1,2,3\}$ has tail exactly $\tfrac1{48}=1/R$
(Lemma~2.1), with $R=|P(1)|=|P_2(1)|=48$ and $\pi=1$.  There the strict
bound comes from the partial sums $1,-16,85,-162,48$ of
$(x-2)(x-3)(x-5)(x-7)$, not from convexity; an exact search gives tail
$0.536/48$ for $X_0=\{1\}$ and $\{2\}$, and $0.634/48$ for $\{3\}$.  So no
counterexample to the verbatim statement exists: the only new equality case
is absorbed by the $(x-2)$ argument of Theorem~\ref{thm:converse}.

\begin{corollary}[When the repaired bound is the infimum, repeated roots]
\label{cor:repairedsharprep}
Corollary~\ref{cor:repairedsharp} holds verbatim for
$1<r_1\le\cdots\le r_L$, repetitions allowed.
\end{corollary}

\begin{proof}
Its proof, with Theorem~\ref{thm:twosidedsharprep} in place of
Theorem~\ref{thm:twosidedsharp}, and Theorem~\ref{thm:converse} with
Corollary~\ref{cor:charrep} in place of Proposition~\ref{prop:partial} with
Theorem~\ref{thm:converse} when $r_1\ge2$.
\end{proof}

When the criterion fails, Theorem~\ref{thm:value} still computes the
infimum as $1/T$: its proof never uses the cutoff.  What does use it is the
covering, through the deletion step in Theorem~\ref{thm:lift}, and below $2$
the escaping placement need not be extremal.  For one exempt coefficient the
covering survives, with the half-mass point of
Corollary~\ref{cor:halfmass}, which needs no cutoff, in place of $\max Z^*$.

The half-mass point comes from a line of certificates, and nothing on that
line needs the cutoff.  Adding one node to a certificate and one zero to its
zero set moves it along a line, on which the tail is a weighted sum of
distances.

\begin{lemma}[Sign rule for the pencil]\label{lem:pencil}
Let $n=L-k-1$, let $A$ be a set of $n-1$ positive integers, let $F$ be the full
certificate on $y_{k+2},\ldots,y_L$ with zero set $A$, and let $H\ne0$ be an
exponential sum on $y_{k+1},\ldots,y_L$ with $H_0=0$ and $H_a=0$ for $a\in A$.
Then $H$ is unique up to scale, the certificates on $y_{k+1},\ldots,y_L$
vanishing on $A$ are exactly $F+tH$ with $t\in\mathbb R$, and for every $d\ge1$
with $d\notin A$
\[
  \sgn F_d=(-1)^{|\{a\in A\,:\,a<d\}|},\qquad
  \sgn H_d=\varepsilon\,(-1)^{1+|\{a\in A\,:\,a<d\}|}
\]
for one fixed $\varepsilon\in\{\pm1\}$.  In particular $F_dH_d$ has a single
sign throughout.
\end{lemma}

\begin{proof}
$H$ solves $n$ homogeneous conditions in $n+1$ unknowns, a system nonsingular
by the zero bound, so it is unique up to scale; adding $tH$ to $F$ preserves
$v_0=1$ and the zeros on $A$, and the dimensions match.  Now $F$ is an
exponential sum on $n$ nodes with the $n-1$ zeros $A$, so by the zero bound it
has no others and changes sign at each of them, and $F_0=1>0$.  Likewise $H$ is
one on $n+1$ nodes with the $n$ zeros $\{0\}\cup A$ and no others; for $d\ge1$
the count of zeros below $d$ includes the one at $0$.
\end{proof}

\begin{proposition}[The displaced zero is a weighted median]\label{prop:median}
In the notation of Lemma~\ref{lem:pencil}, let $N\ge1$ with $N\notin A$.  The
full certificate on $y_{k+1},\ldots,y_L$ with zero set $A\cup\{N\}$ is
$F+t_NH$ with $t_N=-F_N/H_N$, and
\begin{equation}\label{eq:median}
  \Tail\bigl(A\cup\{N\}\bigr)
  =\sum_{d\ge1}\Bigl|\,|F_d|-\theta_N|H_d|\,\Bigr|
  =\sum_{d\ge1}|H_d|\,\Bigl|\theta_N-\frac{|F_d|}{|H_d|}\Bigr|,
  \qquad\theta_N=\frac{|F_N|}{|H_N|}.
\end{equation}
The right-hand side is convex and piecewise linear in $\theta$ with breakpoints
exactly the $\theta_N$, so
\[
  \min_{N\ge1,\;N\notin A}\Tail\bigl(A\cup\{N\}\bigr)
  =\min_{\theta\ge0}\sum_{d\ge1}|H_d|\,\Bigl|\theta-\frac{|F_d|}{|H_d|}\Bigr|,
\]
attained at the weighted median of $\bigl(|F_d|/|H_d|\bigr)_{d\ge1,\,d\notin A}$
with weights $|H_d|$.
\end{proposition}

\begin{proof}
By Lemma~\ref{lem:pencil} the certificates vanishing on $A$ form the line
$F+tH$, and $t_N$ is the unique parameter killing $N$.  There $t_NH$ opposes
$F$, so by the sign rule it opposes $F$ at every $d\ge1$, whence
$|F_d+t_NH_d|=\bigl||F_d|-|t_N|\,|H_d|\bigr|$ with $|t_N|=\theta_N$, which is
\eqref{eq:median}.  Each summand is $|H_d|$ times the modulus of an affine
function of $\theta$ vanishing at $|F_d|/|H_d|$, so the sum is convex and
piecewise linear with those breakpoints and attains its minimum at one of them;
a weighted sum of distances is minimized at the weighted median.
\end{proof}

The median is an order statistic, because the points it ranges over are
already sorted by $d$.

\begin{lemma}[The ratio is monotone]\label{lem:ratio}
In the notation of Lemma~\ref{lem:pencil}, $d\mapsto|F_d|/|H_d|$ is strictly
decreasing on $\{d\ge0:d\notin A\}$.
\end{lemma}

\begin{proof}
For $c\in\mathbb R$ the sequence $H-cF$ is an exponential sum on the $n+1$
nodes $y_{k+1},\ldots,y_L$, and is not identically zero, since $F$ gives weight
$0$ to $y_{k+1}$ and $H$ does not.  By the zero bound it has at most $n$ real
zeros with multiplicity, and it vanishes on $A$, which is $n-1$ of them.  So
$\Phi=H/F$ takes each value at most once off $A$, and takes none of them at an
$a\in A$ as well, as that would be a second zero there.  On each component of
$[0,\infty)\setminus A$ the ratio $\Phi$ is continuous, and $F$ and $H$ have a
simple zero at each $a\in A$, so $\Phi$ extends continuously across it.  Being
continuous and injective on $[0,\infty)$, $\Phi$ is strictly monotone; and
$\Phi(0)=0$ because $H_0=0$, while $|\Phi(d)|\to\infty$ because $|H_d|$ decays
like $y_{k+1}^d$ and $|F_d|$ like $y_{k+2}^d$, with $y_{k+2}<y_{k+1}$.
\end{proof}

\begin{corollary}[The displaced zero is a half-mass point]\label{cor:halfmass}
The minimum in Proposition~\ref{prop:median} is attained at
\[
  N=\max\Bigl\{d\ge1,\ d\notin A\ :\
  \sum_{e\ge d}|H_e|\ \ge\ \tfrac12\sum_{e\ge1}|H_e|\Bigr\}.
\]
\end{corollary}

\begin{proof}
By Lemma~\ref{lem:ratio} the order of the points $|F_d|/|H_d|$ is the reverse
of the order of $d$, so the weights $|H_d|$ of Proposition~\ref{prop:median}
accumulate from the largest $d$ downwards.
\end{proof}

The displaced zero $N$ need not exceed $\max A$.
Taking $\theta=0$ in \eqref{eq:median} returns $\Tail(F)$, and the breakpoints
are positive and tend to $0$, since $|F_d|/|H_d|$ decays like
$(y_{k+2}/y_{k+1})^d$; so the convex minimum is at most the value at $0$.
With nodes anywhere in $(0,1)$, some $N$ therefore gives
$\Tail(A\cup\{N\})\le\Tail(F)$: for a \emph{fixed} zero set the deletion
step reverses beyond the cutoff, as on $(1/3,1/2,9/10)$ in
Section~\ref{sec:extend}, but once the
displaced zero is chosen by Proposition~\ref{prop:median} it cannot.

\begin{proposition}[One exempt coefficient at any roots]\label{prop:uoneall}
Let $1<r_1<\cdots<r_L$ and $u=1$, let $\tau^*$ be attained by the full
certificate $v^*$ on $y_2,\ldots,y_L$ with zero set $Z^*$
(Proposition~\ref{prop:attain}), and let $N'$ be the half-mass point of
Corollary~\ref{cor:halfmass} for $k=0$ and $A=Z^*$, so that $F=v^*$ and $H$ is
the sum on all $L$ nodes vanishing on $\{0\}\cup Z^*$.  Then
\[
  \inf_Qb_2\Bigl(Q\prod_{i=1}^L(x-r_i)\Bigr)
  =\frac1{\max\bigl(\tau^*,\ \max_{w\in G}\tau(\{w\})\bigr)},
  \qquad G=\{1,\ldots,N'-1\}\setminus Z^*.
\]
\end{proposition}

\begin{proof}
By \eqref{eq:median} the tail along the pencil is
$\varphi(\theta)=\sum_{d\ge1}|H_d|\,\bigl|\theta-|F_d|/|H_d|\bigr|$, with
$\varphi(0)=\Tail(v^*)=\tau^*$.  All breakpoints are positive, so the right
derivative at $0$ is $-\sum_{d\notin Z^*}|H_d|<0$, and the convex $\varphi$
has its minimum, at $\theta_{N'}$, below $\tau^*$.  For $0<\theta\le
\theta_{N'}$ convexity gives $\varphi(\theta)\le(1-s)\tau^*+s\varphi(\theta_{N'})
<\tau^*$ with $s=\theta/\theta_{N'}$.  By Lemma~\ref{lem:ratio},
$\theta_N=|F_N|/|H_N|$ is positive and decreasing in $N\notin Z^*$, so
$\Tail(Z^*\cup\{N\})<\tau^*$ for every $N\ge N'$ outside $Z^*$.  That full
certificate is one for $\{w\}$ when $w=N$, and, taking $N=N'$, when
$w\in Z^*$.  So $\tau(\{w\})<\tau^*$ for $w\notin G$, and $T\ge\tau^*$ by
Theorem~\ref{thm:value}.
\end{proof}

Exact values, each $\tau$ certified by a dual multiplier as in
Proposition~\ref{prop:attain} and every $\tau(\{w\})$, $w\in G$, computed:
\[
\begin{array}{lcccc}
  (r_i) & \eqref{eq:repaired} & \eqref{eq:twosided} & \inf b_2 & \text{extremal}\\\hline
  (\tfrac32,10,20) & \tfrac{171}2 & \tfrac{171}2 & \tfrac{171}2 & w=2\\
  (\tfrac32,3,4) & 3 & 3 & 6 & \text{escaping}\\
  (\tfrac32,3,5,7) & 24 & 24 & \tfrac{1704}{31} & \text{escaping}\\
  (\tfrac{11}{10},3,5,7) & \tfrac{24}5 & \tfrac{24}5 & \tfrac{3398808}{96935}=35.06\ldots & w=5\\
  (\tfrac65,\tfrac32,3) & \tfrac15 & \tfrac15 & \tfrac{23688369}{19164415}=1.236\ldots & w=1
\end{array}
\]
At $(\tfrac32,3,4)$ the infimum is the product $(r_2-1)(r_3-1)$ that
\eqref{eq:order} would assert, twice \eqref{eq:repaired}: the exempt root is
absorbed as if it were at least $2$.  At $(\tfrac65,\tfrac32,3)$ the infimum
exceeds even that product, $1$.  At $(\tfrac{11}{10},3,5,7)$ the surviving
roots are those of $(2,3,5,7)$, so $\tau^*=\tfrac{31}{1704}$, but
$\tau(\{5\})=\tfrac{96935}{3398808}>\tau^*$: the finite placement beats the
escaping one, and \eqref{eq:step} fails at $k=0$, $X=\{5\}$.  Below the
cutoff Proposition~\ref{prop:step} can therefore not be applied as is, and
for $u\ge2$ the value $1/T$ remains a search.

\section{Scope and open questions}\label{sec:scope}

\paragraph{What is proved.}  For roots at least $2$ the question whether
\eqref{eq:order} is the infimum is settled, repeated roots included
(Theorem~\ref{thm:converse}, Corollary~\ref{cor:charrep}), and so is the
value when it is not (Theorems~\ref{thm:extremal} and~\ref{thm:extremalrep}).
Below $2$, \eqref{eq:repaired} and \eqref{eq:twosided} hold for every
$r_1>1$, and Theorem~\ref{thm:twosidedsharp}, Corollary~\ref{cor:repairedsharp}
and their repeated-root forms (Theorem~\ref{thm:twosidedsharprep},
Corollary~\ref{cor:repairedsharprep}) say when they are attained.

\paragraph{Open questions.}
\begin{itemize}
\item[(i)] \emph{Below the cutoff with $u\ge2$.}  Theorem~\ref{thm:value} gives the
infimum as $1/T$ for all distinct roots above $1$, but at
$(\tfrac{11}{10},3,5,7)$ the escaping placement is not extremal, so
Proposition~\ref{prop:step} is unavailable.  For $u=1$,
Proposition~\ref{prop:uoneall} reduces $T$ to finitely many placements; for
$u\ge2$ computing $T$ remains a search.
\item[(ii)] \emph{The shift at run ends.}  The hypothesis of
Corollary~\ref{cor:shiftrun} moves the start of each run of $A$ past its end.
The variant that moves the end instead,
$\Tail(B_0\cup A)\le\Tail(B_0\cup A\setminus\{a\}\cup\{a+1\})$ for every
$a\in A$ with $a+1\notin A$, agrees with it when no two elements of $A$ are
consecutive; for longer runs we do not know whether it implies the same
conclusion, though it did in all $550$ instances we tried.  Minimality
supplies both forms, so Theorem~\ref{thm:extremal} does not need it.
\item[(iii)] \emph{Attainment at repeated roots.}  At distinct roots the
strict clause of Theorem~\ref{thm:lift} shows the infima of
Proposition~\ref{prop:values} are not attained.  At repeated roots it is not
carried over, so whether the infima of Proposition~\ref{prop:valuesrep} are
attained is open, though an exact search finds a certificate with tail below
$\tau^*$ for every placement inside $\{1,\ldots,30\}$.
\item[(iv)] \emph{The rate of far zeros.}  Proposition~\ref{prop:farrep}
gives $\eta\to0$ for every shape of $E$ but no rate.  The rate
$O(1/\min E)$ is known at bounded diameter when the cut falls inside a
multiplicity block, and $\eta$ stayed $\Theta(1/\min E)$ in every
configuration computed, including gaps $\lfloor\sqrt n\rfloor$ and geometric
spacing; a rate for arbitrary $E$ is not proved, and no result here needs one.
\end{itemize}

\paragraph{Relation to the literature.}  The bounds here concern coefficient
magnitudes, not ordinary length or height: they are uniform over the degree
of $F$ and over all additional real or complex roots, and
in~\cite{coefficient-mass} the hypothesis $|f_D|\ge1$ enters only where the
$\log|f_D|$ term is discarded to reach an asymptotic.  In the opposite
direction, the degree of an integer polynomial of height $H\ge2$ all of whose
roots are nonzero rationals is at most a constant times $\log H$~\cite{htv},
and the minimal height of a degree-$n$ polynomial over a fixed number field
with all roots in $K^\times$ grows exponentially in $n$~\cite{tromp}; the
bound here constrains neither degree nor height but the additive coefficient
mass.  The contribution of~\cite{coefficient-mass} relative to reduced-height
duality is the prefix-sensitive tail comparison and its iteration; the
normalized determinant sequence and its consecutive-zero evaluation are from
Dubickas and Jankauskas~\cite{dj}, as is the case $u=0$ of the construction
of Section~\ref{sec:criterion}.  Proposition~\ref{prop:confluent} is not needed for
\eqref{eq:order}, which Theorem~3.2 reaches at repeated roots by
perturbation, but Theorem~\ref{thm:converse} uses it at repeated roots.

The infimum of $b_{u+1}$ is the $\ell^\infty$ member of a family of extremal
problems over normalized multiples.  Reduced length minimizes $\ell^1$
over normalized multiples~\cite{schinzel}.  Its integer-polynomial $\ell^2$
analogue, the minimal Euclidean norm of an algebraic number, is effectively
computable~\cite{frw}.  The logarithmic coefficient mass $\Lambda$
of~\cite{coefficient-mass} is additive over coefficients rather than a norm at
all.  Corollary~3.6 settles the asymptotic form of the aggregated question for
prime roots, and the multiplier $x+1$ of~\cite[Section~3]{coefficient-mass}
shows that even there the root product need not be the finite-$L$
minimizer.

\paragraph{Data availability.}
No datasets were generated or analyzed in this study.

\begin{thebibliography}{9}
\small
\bibitem{karlin-studden}
S.~Karlin and W.~J. Studden,
\newblock \emph{Tchebycheff Systems: With Applications in Analysis and Statistics},
\newblock Interscience, New York, 1966.

\bibitem{polya-szego}
G.~P\'olya and G.~Szeg\H{o},
\newblock \emph{Problems and Theorems in Analysis II},
\newblock Springer, 1976, Part Five, Chapter 1,
\newblock \href{https://doi.org/10.1007/978-3-642-61905-2}
{doi:10.1007/978-3-642-61905-2}.

\bibitem{schinzel}
A.~Schinzel,
\newblock On the reduced length of a polynomial with real coefficients,
\newblock \emph{Functiones et Approximatio Commentarii Mathematici} 35 (2006), 271--306,
\newblock \href{https://doi.org/10.7169/facm/1229442629}
{doi:10.7169/facm/1229442629}.

\bibitem{dj}
A.~Dubickas and J.~Jankauskas,
\newblock On the reduced height of a polynomial,
\newblock \emph{Publicationes Mathematicae Debrecen} 71 (2007), 325--348,
\newblock \href{https://doi.org/10.5486/PMD.2007.3728}
{doi:10.5486/PMD.2007.3728}.

\bibitem{frw}
M.~Filaseta, M.~L. Robinson and F.~S. Wheeler,
\newblock The minimal Euclidean norm of an algebraic number is effectively
computable,
\newblock \emph{Journal of Algorithms} 16 (1994), 309--333,
\newblock \href{https://doi.org/10.1006/jagm.1994.1015}
{doi:10.1006/jagm.1994.1015}.

\bibitem{htv}
L.~Hajdu, R.~Tijdeman and N.~Varga,
\newblock On polynomials with only rational roots,
\newblock \emph{Mathematika} 69 (2023), 867--878,
\newblock \href{https://doi.org/10.1112/mtk.12209}
{doi:10.1112/mtk.12209}.

\bibitem{tromp}
T.~Tromp,
\newblock On the height of polynomials that split completely over a fixed
number field,
\newblock arXiv:2509.12064 (2025),
\newblock \href{https://doi.org/10.48550/arXiv.2509.12064}
{doi:10.48550/arXiv.2509.12064}.

\bibitem{coefficient-mass}
B.~Pham,
\newblock Displaced-zero certificates and the coefficient mass of polynomial
multiples,
\newblock manuscript in preparation;
Lemma~2.1 (Consecutive zero set), Theorem~2.2 (Displaced-zero comparison),
Lemma~2.3 (Interlacing of the correction sum), Lemma~3.1 (Slack
certificate), Theorem~3.2 (Coefficient order statistics),
Corollary~3.6 (Sharp prime-root infimum), Proposition~4.5 (Sharpness of the
row $k=2$ at $(2,3)$).
\end{thebibliography}

\end{document}
